\documentclass[11pt]{article}
% Shared preamble for the rigor documents (Lamport-style structured proofs).  \input this file after \documentclass.
\usepackage[T1]{fontenc}\usepackage{lmodern}\usepackage{microtype}
\usepackage[margin=1in]{geometry}
\usepackage{amsmath,amssymb,amsthm,mathtools}
\usepackage{xcolor}\usepackage{enumitem}\usepackage{booktabs}\usepackage{graphicx}\usepackage{hyperref}
\usepackage[most]{tcolorbox}
\definecolor{navy}{HTML}{17365D}\definecolor{teal}{HTML}{117A8B}\definecolor{warn}{HTML}{A33A2B}\definecolor{okgreen}{HTML}{2E7D4F}
\hypersetup{colorlinks=true,linkcolor=navy,citecolor=teal,urlcolor=teal}
\theoremstyle{plain}
\newtheorem{theorem}{Theorem}[section]\newtheorem{lemma}[theorem]{Lemma}\newtheorem{proposition}[theorem]{Proposition}\newtheorem{corollary}[theorem]{Corollary}
\theoremstyle{definition}\newtheorem{definition}[theorem]{Definition}\newtheorem{remark}[theorem]{Remark}\newtheorem{claim}[theorem]{Claim}\newtheorem{issue}[theorem]{Issue}
% ---------- Lamport structured proofs ----------
% \lstep{level}{number}{assertion}   -> indented "<level>number. assertion"
% \lproof                             -> "PROOF:" line (start of the sub-proof of the preceding step)
% \lby{justification}                 -> "BY ..." leaf justification for the preceding step
% \lqed{level}{number}                -> QED step of a sub-proof
% keywords: \lassume \lprove \lsuffices \lcase \ldefine \lpick \lnew
\newlength{\lindent}\setlength{\lindent}{1.7em}
\newcommand{\lstep}[3]{\par\smallskip\noindent\hspace*{\dimexpr\numexpr#1-1\relax\lindent\relax}{\color{navy}$\langle#1\rangle#2$.}\ #3}
\newcommand{\lproof}{\par\noindent\hspace*{\lindent}{\scshape Proof:}}
\newcommand{\lby}[1]{\par\noindent\hspace*{\lindent}{\scshape By}\ #1.}
\newcommand{\lqed}[2]{\par\smallskip\noindent\hspace*{\dimexpr\numexpr#1-1\relax\lindent\relax}{\color{navy}$\langle#1\rangle#2$.}\ {\scshape Q.E.D.}}
\newcommand{\lassume}{{\scshape Assume}\ }\newcommand{\lprove}{{\scshape Prove}\ }\newcommand{\lsuffices}{{\scshape Suffices}\ }
\newcommand{\lcase}{{\scshape Case}\ }\newcommand{\ldefine}{{\scshape Define}\ }\newcommand{\lpick}{{\scshape Pick}\ }\newcommand{\lnew}{{\scshape New}\ }
% ---------- verdict boxes ----------
\newtcolorbox{verdictok}{colback=okgreen!8,colframe=okgreen,title={\bfseries Verdict: verified},fonttitle=\small}
\newtcolorbox{verdictgap}{colback=orange!10,colframe=orange!80!black,title={\bfseries Verdict: gap / caveat},fonttitle=\small}
\newtcolorbox{verdictbreak}{colback=warn!10,colframe=warn,title={\bfseries Verdict: proof-breaking issue},fonttitle=\small}
\newtcolorbox{citedbox}[1]{colback=navy!5,colframe=navy!60,title={\bfseries Cited result: #1},fonttitle=\small,breakable}
\setlength{\parindent}{0pt}\setlength{\parskip}{4pt}\allowdisplaybreaks
\newcommand{\cA}{\mathcal A}\newcommand{\cF}{\mathcal F}\newcommand{\cM}{\mathfrak M}\newcommand{\cN}{\mathfrak N}

% extra calligraphic shortcuts (\providecommand: no clash if a document defines its own)
\providecommand{\cB}{\mathcal B}\providecommand{\cC}{\mathcal C}\providecommand{\cD}{\mathcal D}\providecommand{\cE}{\mathcal E}\providecommand{\cG}{\mathcal G}\providecommand{\cH}{\mathcal H}\providecommand{\cI}{\mathcal I}\providecommand{\cJ}{\mathcal J}\providecommand{\cK}{\mathcal K}\providecommand{\cL}{\mathcal L}\providecommand{\cO}{\mathcal O}\providecommand{\cP}{\mathcal P}\providecommand{\cQ}{\mathcal Q}\providecommand{\cR}{\mathcal R}\providecommand{\cS}{\mathcal S}\providecommand{\cT}{\mathcal T}\providecommand{\cU}{\mathcal U}\providecommand{\cV}{\mathcal V}\providecommand{\cW}{\mathcal W}\providecommand{\cX}{\mathcal X}\providecommand{\cY}{\mathcal Y}\providecommand{\cZ}{\mathcal Z}
\newcommand{\Fid}{\operatorname{F}}\newcommand{\Tr}{\operatorname{Tr}}\newcommand{\eps}{\varepsilon}\newcommand{\dd}{\,\mathrm d}

\providecommand{\one}{\mathbf 1}\providecommand{\CAR}{\mathrm{CAR}}\providecommand{\Ran}{\operatorname{Ran}}\providecommand{\dist}{\operatorname{dist}}\providecommand{\norm}[1]{\lVert #1\rVert}
\title{The ultraviolet tail of the recovery fidelity:\\
a rigorous bound, error bars, and the rate}
\author{Agent PT (Phase 1, problem O4)}
\date{2026-09-08}
\begin{document}
\maketitle
\tableofcontents
\section{Setting, notation, and what is proved}\label{sec:setting}

\paragraph{Conventions.}  We keep the notation of Note~5 (\texttt{fidelity\_recovered\_quasifree.tex}).
$r=1$ fermion component throughout (all quantities are proportional to $r$; restore it by multiplying
$\Phi$, $\tau_\pm$, $g_B$ by $r$).  $I=ABC=(-a,L)$, $A=(-a,0)$, $D=BC=(0,L)$, $H=L^2(I)$,
$Q=E_IP_+E_I$ the vacuum symbol, $\widetilde Q=\widetilde Q_s$ the symbol of the zero-collar recovered
state, $\delta=Q-\widetilde Q$ the defect, $\zeta=as/(a+L)$ the cross ratio variable,
\[
 \Fid(\zeta)=\Fid_{\cF(I)}(\omega,\widetilde\omega_s),\qquad \Phi(\zeta)=-\log\Fid(\zeta).
\]
In the modular ($\kappa$-)representation, $Q$ is the multiplication operator $q_\kappa=(1+e^\kappa)^{-1}$
on $L^2(\mathbb R,\dd\kappa)$ (Note~5, (\texttt{modes})--(\texttt{fermi})); $K=\log\frac{1-Q}{Q}$ is
multiplication by $\kappa$.  For $\kappa_c>0$ put
\[
 E:=\one_{[-\kappa_c,\kappa_c]}(K),\quad E^\perp=1-E,\quad E^+=\one_{(\kappa_c,\infty)}(K),\quad
 E^-=\one_{(-\infty,-\kappa_c)}(K),
\]
so that $\eps:=e^{-\kappa_c}$ is (up to $1+O(\eps)$) the spectral cut of Note~5, \S6:
$\{\eps<Q<1-\eps\}=\one_{(-\kappa_c,\kappa_c)}(K)$ exactly, with $\eps=(1+e^{\kappa_c})^{-1}$.
We write $h(t):=\sqrt{t(1-t)}$ for $t\in[0,1]$, and for a symbol $0\le X\le1$
\[
 \Theta_X:=\begin{pmatrix}X^{1/2}\\ (1-X)^{1/2}\end{pmatrix}:H\to H\oplus H,
 \qquad
 \mathbb X:=\Theta_X\Theta_X^*=\begin{pmatrix}X&h(X)\\ h(X)&1-X\end{pmatrix}.
\]
$\|\cdot\|_1,\|\cdot\|_2,\|\cdot\|$ are the trace, Hilbert--Schmidt and operator norms.

\paragraph{The window quantities.}  $\Phi_W(\zeta):=-\log\Fid\bigl(\omega_{EQE},\omega_{E\widetilde QE}\bigr)$
computed inside $EH$ --- the ``spectral-window compression'' of Note~5, \S6, in the exact continuum
(the box is treated separately in \S\ref{sec:box}).  The two ultraviolet occupation tails of the
recovered state are
\begin{equation}\label{eq:taudef}
 \tau_+:=\Tr\bigl[E^+\widetilde QE^+\bigr]=\int_{\kappa_c}^{\infty}\widetilde Q(\kappa,\kappa)\dd\kappa,
 \qquad
 \tau_-:=\Tr\bigl[E^-(1-\widetilde Q)E^-\bigr],
\end{equation}
and the corresponding vacuum tails $\nu_\pm$, $\nu_+=\int_{\kappa_c}^\infty q_\kappa\dd\kappa
=\log(1+e^{-\kappa_c})\le\eps$, $\nu_-=\nu_+$.  Put $\tau:=\tau_++\tau_-$, $\nu:=\nu_++\nu_-\le2\eps$.

\paragraph{What is proved.}  Three statements, in decreasing order of strength.
\begin{enumerate}[label=(\Alph*),leftmargin=2.2em]
\item \emph{A structural tail bound with no kernel estimates} (Theorem~\ref{thm:main}, \S\ref{sec:tail}):
 \[
  0\ \le\ \Phi(\zeta)-\Phi_W(\zeta)\ \le\ \Psi\bigl(\Phi_W,\ \Phi_\partial\bigr),\qquad
  \Phi_\partial\le\frac{5\tau+\nu}{2\,(1-5\tau-\nu)},
 \]
 with $\Psi$ the explicit triangle-inequality function \eqref{eq:Psi} (it also involves $\Phi$; valid under the hypotheses $5\tau+\nu<1$ and $Y<1$ of Theorem~\ref{thm:main}, with $Y=e^{\Phi_W}y$ in place of the printed $y$, corrected 2026-09-24, see Correction~C1); to leading order
 $\Phi-\Phi_W\le2\sqrt{\Phi_W\,\Phi_\partial}+\Phi_\partial$.  Only positivity of $\widetilde Q$,
 operator concavity of $h$, and the definition of $\tau_\pm$ enter.
\item \emph{An occupation-tail estimate with a controlled remainder}
 (Theorem~\ref{thm:tau}, \S\ref{sec:delta}): for $0<\zeta\le1$ and $\kappa_c\ge8$,
 \[
  \Bigl|\tau_\pm-\nu_\pm-\frac{\zeta^2}{15\,\kappa_c^{2}}\Bigr|
  \ \le\ \frac{C_2\,\zeta^{2}}{\kappa_c^{3}}\ +\ C_1\,\zeta\,(1+\kappa_c)\,e^{-\kappa_c}\qquad(\text{corrected 2026-09-24, see Correction~C7}),
 \]
 with $C_1\le0.051$ and $C_2\le0.47$ (printed as $C_2\le0.46$ before 2026-09-24, see Correction~C10; both evaluated numerically in double precision in \S\ref{sec:num}, not certified; printed as ``certified numerically'' before 2026-09-24, see Correction~C4; the analytic
 argument gives the shape and finiteness of the constants, their numerical values are evaluated).
\item \emph{Consequences} (\S\ref{sec:tail}, \S\ref{sec:rate}): $\Phi-\Phi_W\le C\,\zeta^2/\kappa_c$
 with a finite $C$ on the region of Corollary~\ref{cor:rate} (the explicit $C=0.173$ there, with a further $0.88\,\zeta^2/\kappa_c^2$, rests on numerical inputs; printed as ``with an explicit $C$'' before 2026-09-24, see Correction~C4) --- one power of $\kappa_c$ short of the conjectured
 $\zeta^2/(15\kappa_c^2)$, the loss being the cross term of the Bures triangle inequality; and,
 \emph{conditionally on} the finite-dimensional cubic estimate of \texttt{quadratic\_limit.tex},
 Problem~7.2, the rate $\Phi=f_2\zeta^2\bigl[1+O(\log^{-1}(1/\zeta))\bigr]$ --- one logarithm short of
 the $\log^{-2}$ rate claimed by an earlier version of Note~5, Cor.~4.3 (printed as ``Cor.~4.2'' before 2026-09-24; the current version says that this rate was wrong, see Correction~C9).
\end{enumerate}
Sections~\ref{sec:box}--\ref{sec:bars} turn (A)+(B) into error bars for the tabulated values (not certified: their upper half rests on the numerical inputs of \eqref{eq:final}; printed as ``certified error bars'' before 2026-09-24, see Correction~C8),
and \S\ref{sec:num} verifies every inequality on the free-fermion data.
\section{The window state and the exact factorisation of the fidelity}\label{sec:window}

\paragraph{Standing finite-dimensional convention.}  All determinant and trace identities below are
first proved on a finite-dimensional subspace and then transported to $\cF(I)$ by the monotone limit
of Note~5, Theorem~3.1.  Fix once and for all an increasing sequence of finite-rank projections
$P_n\uparrow1$ on $H\cong L^2(\mathbb R,\dd\kappa)$ with $[P_n,E]=0$ (take an orthonormal basis of
$L^2([-\kappa_c,\kappa_c])$ and one of $L^2(\{|\kappa|>\kappa_c\})$ and let $P_n$ project onto the
first $n$ vectors of each).  Every symbol $X$ below is replaced by $P_nXP_n$ on $P_nH$; all bounds
proved are uniform in $n$, and $\Fid$, $\Phi$, $\Phi_W$ are the (monotone) limits.  We suppress the
index $n$.

\begin{definition}\label{def:Qpp}
$\displaystyle Q'':=Q-E\delta E=\bigl(E\widetilde QE\bigr)\oplus\bigl(E^\perp QE^\perp\bigr)$,
and $\rho'':=\omega_{Q''}$, $\rho:=\omega_Q$, $\sigma:=\omega_{\widetilde Q}$.
\end{definition}

\begin{lemma}[the window state]\label{lem:factor}
$0\le Q''\le1$ with $\ker Q''=\ker(1-Q'')=\{0\}$, $Q''$ commutes with $E$, and
\[
 \Fid(\rho,\rho'')=\Fid\bigl(\omega_{EQE},\omega_{E\widetilde QE}\bigr)=e^{-\Phi_W}.
\]
\end{lemma}
\lstep{1}{1}{$EQE^\perp=0$, hence $Q=(EQE)\oplus(E^\perp QE^\perp)$ and
$Q-E\delta E=(EQE-E\delta E)\oplus(E^\perp QE^\perp)=(E\widetilde QE)\oplus(E^\perp QE^\perp)$.}
\lby{$E$ is a spectral projection of $K=\log\frac{1-Q}Q$, so $[E,Q]=0$; and
$\delta=Q-\widetilde Q$ by (\texttt{defect-recall})}
\lstep{1}{2}{$0\le E\widetilde QE\le1$ and $0\le E^\perp QE^\perp\le1$, with trivial kernels and
co-kernels on $EH$, $E^\perp H$ respectively.}
\lby{$0\le\widetilde Q\le1$ and $0\le Q\le1$, so their compressions are between $0$ and $1$;
$\langle f,\widetilde Qf\rangle=0$ forces $\widetilde Q^{1/2}f=0$ and $\ker\widetilde Q=\{0\}$
(Note~5, \S3), same for $1-\widetilde Q$ and for $Q$}
\lstep{1}{3}{Under the canonical isomorphism $\Lambda(EH\oplus E^\perp H)\cong\Lambda(EH)\otimes
\Lambda(E^\perp H)$ one has $\rho=\omega_{EQE}\otimes\omega_{E^\perp QE^\perp}$ and
$\rho''=\omega_{E\widetilde QE}\otimes\omega_{E^\perp QE^\perp}$.}
\lby{$\rho_X=\det(1-X)\,\Gamma(G_X)$ with $G_X=X(1-X)^{-1}$ (Note~5, (\texttt{rho-Q})), and
$\Gamma(G_1\oplus G_2)=\Gamma(G_1)\otimes\Gamma(G_2)$, $\det(1-X_1\oplus X_2)=\det(1-X_1)\det(1-X_2)$,
by (\texttt{fock-functor}); both $Q$ and $Q''$ are block diagonal by $\langle1\rangle1$}
\lstep{1}{4}{$\Fid(\rho,\rho'')=\Fid(\omega_{EQE},\omega_{E\widetilde QE})\cdot
\Fid(\omega_{E^\perp QE^\perp},\omega_{E^\perp QE^\perp})=\Fid(\omega_{EQE},\omega_{E\widetilde QE})$.}
\lby{multiplicativity of the root fidelity on tensor products,
$\Fid(\rho_1\otimes\rho_2,\sigma_1\otimes\sigma_2)=\Fid(\rho_1,\sigma_1)\Fid(\rho_2,\sigma_2)$
(immediate from $\Fid=\Tr|\sqrt{\rho}\sqrt\sigma|$ and $\sqrt{\rho_1\otimes\rho_2}=
\sqrt{\rho_1}\otimes\sqrt{\rho_2}$), together with $\Fid(\varphi,\varphi)=1$}
\lqed{1}{5}\lby{$\langle1\rangle1$--$\langle1\rangle4$ and the definition of $\Phi_W$}

\begin{remark}
This is the precise sense in which the numerically computed ``sub-compressed'' value is a fidelity of
a genuine pair of states on the whole algebra and not merely a compression: $\Phi_W=-\log\Fid(\rho,\rho'')$,
with the \emph{same} $\rho$ on the left.  In particular $\Phi_W\le\Phi$ is not needed as an input; it
follows from Theorem~\ref{thm:main} below, and independently from monotonicity of $\Fid$ under the
conditional expectation onto $\cF(EH)$.
\end{remark}
\section{Bures angle, and a purification lower bound for the fidelity}\label{sec:triangle}

\begin{citedbox}{Bures angle is a metric (Nielsen--Chuang 2000, Thm.~9.9 / Bengtsson--\.Zyczkowski 2006, \S9.4)}
For states $\rho,\sigma$ put $A(\rho,\sigma):=\arccos\Fid(\rho,\sigma)\in[0,\pi/2]$.  Then $A$ is a
metric on the state space of any von Neumann algebra; in particular the triangle inequality
$A(\rho,\sigma)\le A(\rho,\rho'')+A(\rho'',\sigma)$ holds.  (For matrix algebras this is
Nielsen--Chuang, Theorem~9.9, p.~413, ``the angle $\arccos F$ is a metric''; the general case follows
because $\Fid$ on a von Neumann algebra is the infimum of the fidelities of the finite-dimensional
restrictions and $\arccos$ is monotone --- or directly from Alberti--Uhlmann's transition probability.
\emph{Usage here}: with $\rho''$ of Definition~\ref{def:Qpp}.  \emph{Verdict}: hypotheses (normal
states on a von Neumann algebra) are satisfied; we use it only through the finite-dimensional
compressions, where it is the matrix statement.)
\end{citedbox}

\begin{lemma}[purification lower bound]\label{lem:purif}
Let $0\le X,Y\le1$ on a finite-dimensional $\mathfrak h$ and put
\[
 N(X,Y):=X^{1/2}(1-Y)^{1/2}-(1-X)^{1/2}Y^{1/2}.
\]
Then $\|N(X,Y)\|\le1$ and
\begin{equation}\label{eq:purif}
 \Fid(\omega_X,\omega_Y)\ \ge\ \bigl|\det\bigl(X^{1/2}Y^{1/2}+(1-X)^{1/2}(1-Y)^{1/2}\bigr)\bigr|
 =\det\bigl(1-N^*N\bigr)^{1/2},
\end{equation}
and consequently, whenever $\|N\|_2<1$,
\begin{equation}\label{eq:purif2}
 -\log\Fid(\omega_X,\omega_Y)\ \le\ -\tfrac12\log\det(1-N^*N)\ \le\ \frac{\|N\|_2^2}{2\bigl(1-\|N\|_2^2\bigr)}.
\end{equation}
Moreover
\begin{equation}\label{eq:Nsplit}
 \|N(X,Y)\|_2^2=\tfrac12\|\mathbb X-\mathbb Y\|_2^2=\|X-Y\|_2^2+\|h(X)-h(Y)\|_2^2 .
\end{equation}
\end{lemma}
\lstep{1}{1}{$\Theta_X$ is an isometry $\mathfrak h\to\mathfrak h\oplus\mathfrak h$ with
$\Theta_X\Theta_X^*=\mathbb X$ a projection of rank $n=\dim\mathfrak h$, and
$\Theta_X^\perp:=\binom{-(1-X)^{1/2}}{X^{1/2}}$ is an isometry with
$\Theta_X^\perp\Theta_X^{\perp*}=1-\mathbb X$.}
\lby{$\Theta_X^*\Theta_X=X+(1-X)=1$; $\Theta_X^{\perp*}\Theta_X^\perp=1$;
$\Theta_X^{\perp*}\Theta_X=-(1-X)^{1/2}X^{1/2}+X^{1/2}(1-X)^{1/2}=0$ (both factors are functions of $X$,
hence commute); so the two ranges are orthogonal and together span $\mathfrak h\oplus\mathfrak h$}
\lstep{1}{2}{Let $\Psi_X\in\Lambda(\mathfrak h\oplus\mathfrak h)$ be a unit vector spanning
$\Lambda^{n}(\Ran\mathbb X)$ (a Slater determinant).  Then $\Psi_X$ is the pure gauge-invariant
quasi-free state with symbol $\mathbb X$, and $\langle\Psi_X,\Psi_Y\rangle=\det(\Theta_X^*\Theta_Y)$
up to a phase.}
\lby{a Slater determinant $a^*(f_1)\cdots a^*(f_n)|0\rangle$ has two-point function
$\langle a^*(g)a(f)\rangle=\langle f,\mathbb Xg\rangle$ with $\mathbb X$ the projection onto
$\mathrm{span}(f_i)$; and $\langle a^*(f_1)\cdots a^*(f_n)\Omega,a^*(g_1)\cdots a^*(g_n)\Omega\rangle
=\det[\langle f_i,g_j\rangle]$, applied to $f_i=\Theta_Xe_i$, $g_j=\Theta_Ye_j$ for an orthonormal
basis $(e_i)$ of $\mathfrak h$}
\lstep{1}{3}{The restriction of the state $\Psi_X$ to $\CAR(\mathfrak h\oplus0)$ is $\omega_X$.}
\lby{the restriction of a gauge-invariant quasi-free state to the CAR algebra of a subspace is the
quasi-free state with the compressed symbol (Note~5, proof of Theorem~3.1), and
$\mathbb X_{11}=X$}
\lstep{1}{4}{$\Fid(\omega_X,\omega_Y)\ge|\langle\Psi_X,\Psi_Y\rangle|$.}
\lby{$\langle1\rangle3$ and monotonicity of the root fidelity under the (unital, completely positive)
restriction map, together with $\Fid=|\langle\cdot,\cdot\rangle|$ for two pure states}
\lstep{1}{5}{$|\det(\Theta_X^*\Theta_Y)|^2=\det(\Theta_Y^*\mathbb X\Theta_Y)
=\det(1-\Theta_Y^*(1-\mathbb X)\Theta_Y)=\det(1-N^*N)$, and
$\Theta_X^*\Theta_Y=X^{1/2}Y^{1/2}+(1-X)^{1/2}(1-Y)^{1/2}$,
$\Theta_X^{\perp*}\Theta_Y=N(X,Y)$.}
\lby{$\det(M^*M)=|\det M|^2$; $\Theta_Y^*\Theta_Y=1$; $\langle1\rangle1$; matrix multiplication}
\lstep{1}{6}{$-\tfrac12\log\det(1-N^*N)=-\tfrac12\sum_j\log(1-n_j^2)\le
\tfrac12\frac{\sum_jn_j^2}{1-\max_jn_j^2}\le\frac{\|N\|_2^2}{2(1-\|N\|_2^2)}$, where $n_j$ are the
singular values of $N$.}
\lby{$-\log(1-x)\le x/(1-x)$ for $0\le x<1$, applied with $x=n_j^2\le\max_jn_j^2=\|N\|^2\le\|N\|_2^2$}
\lstep{1}{7}{$\|N\|_2^2=\Tr[\Theta_Y^*(1-\mathbb X)\Theta_Y]=\Tr[(1-\mathbb X)\mathbb Y]
=\tfrac12\|\mathbb X-\mathbb Y\|_2^2=\|X-Y\|_2^2+\|h(X)-h(Y)\|_2^2$.}
\lby{$\Tr[\mathbb X]=\Tr[\mathbb Y]=n$, hence
$\|\mathbb X-\mathbb Y\|_2^2=\Tr[\mathbb X+\mathbb Y-2\mathbb X\mathbb Y]=2\Tr[(1-\mathbb X)\mathbb Y]$;
and the block form of $\mathbb X-\mathbb Y$ has diagonal blocks $\pm(X-Y)$ and off-diagonal blocks
$h(X)-h(Y)$}
\lqed{1}{8}\lby{$\langle1\rangle4$--$\langle1\rangle7$}

\begin{remark}[sharpness]
\eqref{eq:purif} is the ``canonical purification'' (thermofield-double) bound; by Note~5,
Theorem~2.1 it is an equality exactly when the unitary $U$ of the polar decomposition of
$G_X^{1/2}G_Y^{1/2}$ is $1$, in particular whenever $[X,Y]=0$.  It is the finite-dimensional analogue
of the Bogoliubov-extension bound of \texttt{uhlmann\_upper\_bound.tex}; here it is applied not to
$(\rho,\sigma)$ --- for which it is not sharp at second order --- but to the pair $(\rho'',\sigma)$,
whose symbols differ only in the ultraviolet, where they \emph{almost} commute.
\end{remark}
\section{The discarded block: reduction to the occupation tails}\label{sec:delta}

Throughout this section $X:=\widetilde Q$, $Y:=Q''$ of Definition~\ref{def:Qpp}, and we use the block
notation $X_{11}=EXE$, $X_{12}=EXE^\perp$, $X_{22}=E^\perp XE^\perp$, $Q_{22}=E^\perp QE^\perp$, all on
the finite-dimensional space $P_nH$ of \S\ref{sec:window} (so $X_{11}>0$ is invertible).
Recall $\tau_\pm$, $\nu_\pm$ from \eqref{eq:taudef} and $\tau=\tau_++\tau_-$, $\nu=\nu_++\nu_-$.

\begin{theorem}[the discarded block is controlled by the occupation tails]\label{thm:block}
\begin{equation}\label{eq:block}
 \bigl\|N(\widetilde Q,Q'')\bigr\|_2^{2}\ \le\ 5\,\tau+\nu .
\end{equation}
No property of the defect kernel is used: only $0\le\widetilde Q\le1$, $[Q,E]=0$ and the definitions.
\end{theorem}
\lstep{1}{1}{\ldefine $S:=h(X_{11})$, $R:=Eh(X)E|_{EH}$, $R':=E^\perp h(X)E^\perp|_{E^\perp H}$.
Then
\[
 \|N(X,Y)\|_2^2=\underbrace{2\Tr[S^2]-2\Tr[SR]}_{(\mathrm I)}
 +\underbrace{\Tr\bigl[X_{22}(1-Q_{22})\bigr]+\Tr\bigl[Q_{22}(1-X_{22})\bigr]-2\Tr[R'h(Q_{22})]}_{(\mathrm{II})}.
\]}
\lproof
\lstep{2}{1}{$\|N(X,Y)\|_2^2=\Tr[X+Y-2XY]-2\Tr[h(X)h(Y)]$.}
\lby{\eqref{eq:Nsplit}: $\|X-Y\|_2^2+\|h(X)-h(Y)\|_2^2
=\Tr[X^2+Y^2-2XY]+\Tr[h(X)^2+h(Y)^2-2h(X)h(Y)]$ and $h(Z)^2=Z-Z^2$ for $0\le Z\le1$}
\lstep{2}{2}{$Y=X_{11}\oplus Q_{22}$, hence $h(Y)=h(X_{11})\oplus h(Q_{22})=S\oplus h(Q_{22})$,
$\Tr[XY]=\Tr[X_{11}^2]+\Tr[X_{22}Q_{22}]$ and $\Tr[h(X)h(Y)]=\Tr[RS]+\Tr[R'h(Q_{22})]$.}
\lby{Lemma~\ref{lem:factor}$\langle1\rangle1$ and Definition~\ref{def:Qpp}; functional calculus
respects a direct sum; $\Tr[AB]=\sum_{\alpha}\Tr[E^\alpha AE^\alpha B]$ when $B$ is block diagonal
($E^\alpha\in\{E,E^\perp\}$)}
\lstep{2}{3}{$\Tr[X]+\Tr[Y]-2\Tr[XY]
=2\bigl(\Tr X_{11}-\Tr X_{11}^2\bigr)+\Tr[X_{22}(1-Q_{22})]+\Tr[Q_{22}(1-X_{22})]$ and
$\Tr X_{11}-\Tr X_{11}^2=\Tr[S^2]$.}
\lby{$\Tr X=\Tr X_{11}+\Tr X_{22}$, $\Tr Y=\Tr X_{11}+\Tr Q_{22}$, $\langle2\rangle2$, and
$S^2=X_{11}(1-X_{11})$}
\lqed{2}{4}\lby{$\langle2\rangle1$--$\langle2\rangle3$}
\lstep{1}{2}{$0\le R\le S$.}
\lby{$h(t)=\sqrt{t(1-t)}=t\,\#\,(1-t)$ is operator concave on $[0,1]$ (the geometric mean
$(a,b)\mapsto a\#b$ is jointly operator concave, Kubo--Ando 1980, Thm.~3.5; equivalently
$h$ is a nonnegative operator monotone-type function of the form $\sqrt{t}\sqrt{1-t}$, and
operator concavity of $t\mapsto\sqrt{t(1-t)}$ on $[0,1]$ is classical), so Jensen's operator
inequality for an isometry (Hansen--Pedersen 2003, Thm.~2.1) gives $V^*h(X)V\le h(V^*XV)$ with
$V:EH\hookrightarrow H$; and $h\ge0$ gives $R\ge0$}
\lstep{1}{3}{$(\mathrm I)\le2\bigl(\Tr[S^2]-\Tr[R^2]\bigr)$.}
\lby{$\Tr[SR]-\Tr[R^2]=\Tr[(S-R)R]\ge0$, the trace of a product of two positive semi-definite
operators being nonnegative, using $\langle1\rangle2$}
\lstep{1}{4}{$\Tr[S^2]-\Tr[R^2]=\|X_{12}\|_2^2+\|E^\perp h(X)E\|_2^2$.}
\lproof
\lstep{2}{1}{$\Tr[R^2]=\Tr[Eh(X)^2E]-\|E^\perp h(X)E\|_2^2$.}
\lby{$\Tr[(Eh(X)E)^2]=\Tr[Eh(X)(E+E^\perp)h(X)E]-\Tr[Eh(X)E^\perp h(X)E]$}
\lstep{2}{2}{$\Tr[Eh(X)^2E]=\Tr[EX(1-X)E]=\Tr X_{11}-\Tr X_{11}^2-\|X_{12}\|_2^2=\Tr[S^2]-\|X_{12}\|_2^2$.}
\lby{$h(X)^2=X-X^2$ and $(X^2)_{11}=X_{11}^2+X_{12}X_{12}^*$}
\lqed{2}{3}\lby{$\langle2\rangle1$, $\langle2\rangle2$}
\lstep{1}{5}{$\|E^\perp h(X)E\|_2^2\le\Tr[E^\perp X(1-X)E^\perp]\le\tau_++\tau_-=\tau$.}
\lby{$\|E^\perp h(X)E\|_2^2\le\|E^\perp h(X)\|_2^2=\Tr[E^\perp h(X)^2E^\perp]$; then
$0\le X(1-X)\le X$ and $0\le X(1-X)\le1-X$, applied on $E^+$ and on $E^-$ respectively, and
$\Tr[E^+XE^+]=\tau_+$, $\Tr[E^-(1-X)E^-]=\tau_-$}
\lstep{1}{6}{$\|X_{12}\|_2^2\le\tau$.}
\lproof
\lstep{2}{1}{$\|EXE^+\|_2^2\le\|X_{11}\|\cdot\Tr[E^+XE^+]\le\tau_+$.}
\lby{$X\ge0$ with $X_{11}>0$ gives the Schur complement inequality
$(EXE^+)^*X_{11}^{-1}(EXE^+)\le E^+XE^+$; $X_{11}\le\|X_{11}\|$ gives
$X_{11}^{-1}\ge\|X_{11}\|^{-1}$, hence $(EXE^+)^*(EXE^+)\le\|X_{11}\|(EXE^+)^*X_{11}^{-1}(EXE^+)$;
take traces and use $\|X_{11}\|\le1$}
\lstep{2}{2}{$\|EXE^-\|_2^2=\|E(1-X)E^-\|_2^2\le\|(1-X)_{11}\|\Tr[E^-(1-X)E^-]\le\tau_-$.}
\lby{$EXE^-=-E(1-X)E^-$ since $E\one E^-=0$; then $\langle2\rangle1$ with $X$ replaced by $1-X\ge0$}
\lqed{2}{3}\lby{$\|X_{12}\|_2^2=\|EXE^+\|_2^2+\|EXE^-\|_2^2$ and $\langle2\rangle1$, $\langle2\rangle2$}
\lstep{1}{7}{$(\mathrm{II})\le\tau+\nu$.}
\lby{$-2\Tr[R'h(Q_{22})]\le0$ ($R'\ge0$, $h(Q_{22})\ge0$);
$\Tr[X_{22}(1-Q_{22})]\le\Tr[E^+XE^+]+\Tr[E^-(1-Q)E^-]=\tau_++\nu_-$ and
$\Tr[Q_{22}(1-X_{22})]\le\Tr[E^+QE^+]+\Tr[E^-(1-X)E^-]=\nu_++\tau_-$, using
$\Tr[AB]\le\|B\|\Tr A$ for $A\ge0$, $B\ge0$ on each of $E^+$, $E^-$ separately ($Q_{22}$ commutes
with $E^\pm$)}
\lqed{1}{8}\lby{$(\mathrm I)+(\mathrm{II})\le2(\|X_{12}\|_2^2+\|E^\perp h(X)E\|_2^2)+\tau+\nu
\le2(\tau+\tau)+\tau+\nu=5\tau+\nu$, by $\langle1\rangle1$, $\langle1\rangle3$--$\langle1\rangle7$}

\begin{verdictok}
\eqref{eq:block} is proved from positivity alone.  Its content is that the entire ultraviolet error of
the window compression --- including the infrared--ultraviolet off-diagonal couplings $E\delta E^\perp$,
which decay only like $\kappa^{-2}$ and were the source of the ``trap'' remark of Note~5, \S3.1 --- is
bounded by the ultraviolet \emph{occupation} $\tau$ of the recovered state, which is exactly the
quantity Note~5's Proposition~4.1 computes.  The constant $5$ is not optimal (numerically the true
ratio $\|N\|_2^2/(5\tau+\nu)$ is $\approx0.1$, \S\ref{sec:num}).
\end{verdictok}
\subsection{The occupation tail of the recovered state}\label{sec:occ}

We now bound $\tau_\pm$.  Recall from \texttt{check\_note5\_theory.tex}, Prop.~6.1 (closed form) and
\S6.2, the exact representation of the diagonal of the defect in the modular basis: with
$S_p=\sinh\frac p2$, $M=\frac{p+p'}2$,
\begin{equation}\label{eq:Fclosed}
 \widetilde F(p,p')=\frac{\zeta S_pS_{p'}}{\sinh M\,\bigl(\sinh M+2\zeta S_pS_{p'}\bigr)},\qquad
 \widetilde Q(\kappa,\kappa)-q_\kappa=\frac1\pi\operatorname{Im}J(\kappa,\kappa),
\end{equation}
\begin{equation}\label{eq:Jdiag}
 J(\kappa,\kappa)=\frac1{(2\pi)^2}\int_0^\infty e^{i\kappa t/2\pi}\,t\,g(t)\dd t,
 \qquad
 g(t):=\int_0^1\widetilde F\bigl(t\sigma,t(1-\sigma)\bigr)\dd\sigma\ \in\mathbb R .
\end{equation}
Split $g=\zeta g^{(1)}+g^{(2)}$ along the expansion of \eqref{eq:Fclosed} in its own denominator:
\begin{equation}\label{eq:gsplit}
 g^{(1)}(t)=\frac{t\cosh\frac t2-2\sinh\frac t2}{2t\sinh^2\frac t2},
 \qquad
 g^{(2)}(t)=-\int_0^1\frac{2\zeta^2S_{t\sigma}^2S_{t(1-\sigma)}^2}
 {\sinh^2\frac t2\bigl(\sinh\frac t2+2\zeta S_{t\sigma}S_{t(1-\sigma)}\bigr)}\dd\sigma .
\end{equation}

\begin{theorem}[occupation tail with a controlled remainder]\label{thm:tau}
For every $\zeta\in(0,1]$ and $\kappa_c>0$,
\begin{equation}\label{eq:phsym}
 \tau_+=\tau_-=\nu_++\Sigma(\kappa_c),\qquad
 \Sigma(\kappa_c):=\int_{\kappa_c}^\infty\bigl[\widetilde Q(\kappa,\kappa)-q_\kappa\bigr]\dd\kappa,
\end{equation}
and
\begin{equation}\label{eq:taubound}
 \Bigl|\Sigma(\kappa_c)-\frac{\zeta^2}{15\,\kappa_c^{2}}\Bigr|
 \ \le\ \frac{C_2\,\zeta^{2}}{\kappa_c^{3}}\ +\ C_1\,\zeta\,(1+\kappa_c)\,e^{-\kappa_c},
 \qquad C_2=4\pi A_3,\quad C_1=\frac{A_1}{4\pi^2},
\end{equation}
where $A_3:=\sup_{0<\zeta\le1}\zeta^{-2}\|(g^{(2)})'''\|_{L^1(0,\infty)}$ and
$A_1:=\sup_{\kappa_c>0}\bigl|\widehat{g^{(1)}}(\kappa_c/2\pi)\bigr|\,(1+\kappa_c)^{-1}e^{\kappa_c}$
are finite; numerically $A_3\le0.0374$ (printed as $0.036$ before 2026-09-24, see Correction~C10) and $A_1\le2.0$, i.e.
\[
 \boxed{\ C_2\le0.47,\qquad C_1\le0.051\ }\qquad(\text{$C_2\le0.46$ before 2026-09-24, see Correction~C10})
\]
(\S\ref{sec:num}).  Consequently, with $\nu_+=\log(1+e^{-\kappa_c})\le e^{-\kappa_c}$ (the constant $4.7=10\,C_2$ below was printed as $4.6$ before 2026-09-24, see Correction~C10),
\begin{equation}\label{eq:taunu}
 5\tau+\nu\ =\ 12\,\nu_++10\,\Sigma(\kappa_c)
 \ \le\ \frac{2\zeta^2}{3\kappa_c^2}+\frac{4.7\,\zeta^2}{\kappa_c^3}
 +\bigl(12+0.51\,\zeta(1+\kappa_c)\bigr)e^{-\kappa_c}.
\end{equation}
\end{theorem}
\lstep{1}{1}{$g$ is real-valued, real-analytic on $[0,\infty)$ with $g(0)=\frac\zeta6$, and
$|g(t)|+|g'(t)|+|g''(t)|+|g'''(t)|\le C_\zeta e^{-t/2}$ for $t\ge1$.}
\lby{$\widetilde F$ in \eqref{eq:Fclosed} is real and, for $p,p'\ge0$ and $\zeta>0$, its denominator is
$\ge\sinh^2M>0$; the integrand of \eqref{eq:Jdiag} is jointly analytic in $(t,\sigma)$ for
$|t|<2\pi$ (the only zeros of $\sinh\frac t2$ in that disc are at $t=0$, and there numerator and
denominator vanish to the same order, see $\langle1\rangle5$), so $g$ is analytic near $0$;
for large $t$, $S_{t\sigma}S_{t(1-\sigma)}/\sinh^2\frac t2=O(e^{-t/2})$ uniformly in $\sigma$;
differentiation under the $\sigma$-integral is legitimate by dominated convergence}
\lstep{1}{2}{$\displaystyle\Sigma(\kappa_c)=\frac1{2\pi^2}\operatorname{Re}\int_0^\infty g(t)\,
e^{i\lambda t}\dd t$, $\lambda:=\frac{\kappa_c}{2\pi}$.}
\lproof
\lstep{2}{1}{For $\eta>0$, $\int_{\kappa_c}^\infty e^{-\eta\kappa}e^{i\kappa t/2\pi}\dd\kappa
=\dfrac{e^{-\eta\kappa_c}e^{i\kappa_ct/2\pi}}{\eta-it/2\pi}\xrightarrow[\eta\downarrow0]{}
\dfrac{2\pi i\,e^{i\kappa_ct/2\pi}}{t}$ for each $t>0$, with modulus $\le2\pi/t$ uniformly in $\eta$.}
\lby{elementary integration; $|\eta-it/2\pi|\ge t/2\pi$}
\lstep{2}{2}{$\displaystyle\lim_{\eta\downarrow0}\int_{\kappa_c}^\infty e^{-\eta\kappa}J(\kappa,\kappa)\dd\kappa
=\frac1{(2\pi)^2}\int_0^\infty\frac{2\pi i\,e^{i\kappa_ct/2\pi}}{t}\,t\,g(t)\dd t
=\frac i{2\pi}\int_0^\infty g(t)e^{i\lambda t}\dd t$.}
\lby{Fubini (the double integral is absolutely convergent for $\eta>0$ because $|tg(t)|\le C_\zeta te^{-t/2}$
near $\infty$ and $|tg(t)|\le C$ near $0$), then dominated convergence with the majorant
$2\pi|g(t)|\in L^1$ from $\langle2\rangle1$ and $\langle1\rangle1$}
\lstep{2}{3}{$\Sigma(\kappa_c)=\frac1\pi\lim_{\eta\downarrow0}\int_{\kappa_c}^\infty e^{-\eta\kappa}
\operatorname{Im}J(\kappa,\kappa)\dd\kappa
=\frac1\pi\operatorname{Im}\Bigl[\frac i{2\pi}\int_0^\infty ge^{i\lambda t}\dd t\Bigr]
=\frac1{2\pi^2}\operatorname{Re}\int_0^\infty ge^{i\lambda t}\dd t$.}
\lby{\eqref{eq:Fclosed}; $\kappa\mapsto\widetilde Q(\kappa,\kappa)-q_\kappa\ge-q_\kappa$ is integrable on
$(\kappa_c,\infty)$ --- $\widetilde Q(\kappa,\kappa)\ge0$ and
$\int^\infty\widetilde Q(\kappa,\kappa)\dd\kappa<\infty$ because $\widetilde Q-Q$ is trace class
(Note~4: $\|\delta_s\|_1=\frac1{2\pi}\log(1+\zeta)$) --- so monotone/dominated convergence applies
to the Abel limit; and $\operatorname{Im}(iz)=\operatorname{Re}z$ for real $\lambda$-integrals of a
real $g$}
\lqed{2}{4}\lby{$\langle2\rangle1$--$\langle2\rangle3$}

\lstep{1}{3}{$\widetilde Q(-\kappa,-\kappa)-q_{-\kappa}=-\bigl[\widetilde Q(\kappa,\kappa)-q_\kappa\bigr]$,
hence $\tau_-=\tau_+$ and $\tau_+=\nu_++\Sigma(\kappa_c)$, which is \eqref{eq:phsym}.}
\lby{$g$ real and \eqref{eq:Jdiag} give $\operatorname{Im}J(-\kappa,-\kappa)=-\operatorname{Im}J(\kappa,\kappa)$;
therefore $1-\widetilde Q(-\kappa,-\kappa)=1-q_{-\kappa}+[\widetilde Q(\kappa,\kappa)-q_\kappa]
=q_\kappa+[\widetilde Q(\kappa,\kappa)-q_\kappa]=\widetilde Q(\kappa,\kappa)$; integrate over
$\kappa>\kappa_c$ and use $\nu_+=\int_{\kappa_c}^\infty q_\kappa\dd\kappa$}
\lstep{1}{4}{(First order.)  $g^{(1)}$ is even, and
$\operatorname{Re}\int_0^\infty g^{(1)}e^{i\lambda t}\dd t=\tfrac12\widehat{g^{(1)}}(\lambda)$ with
$\widehat{g^{(1)}}(\lambda)=\int_{\mathbb R}g^{(1)}(t)e^{i\lambda t}\dd t$; moreover $g^{(1)}$ is
meromorphic on $\mathbb C$ with double poles exactly at $t\in2\pi i\mathbb Z\setminus\{0\}$ and
$O(e^{-|\operatorname{Re}t|/2})$ decay in the strip $|\operatorname{Im}t|\le2\pi-\eta$, so
$|\widehat{g^{(1)}}(\lambda)|\le A_1(1+2\pi\lambda)e^{-2\pi\lambda}$ with $A_1<\infty$.}
\lby{\texttt{check\_note5\_theory.tex}, Prop.~6.4: $P^{(1)}(-t)=tg^{(1)}(t)\zeta$ is odd with double
poles at $2\pi i\mathbb Z\setminus\{0\}$; hence $g^{(1)}=P^{(1)}/(\zeta t)$ is even and has the same
poles (the simple pole of $1/t$ cancels the triple zero of the numerator at $t=0$, cf.\
$\langle1\rangle5$); $\operatorname{Re}\int_0^\infty g^{(1)}e^{i\lambda t}=\int_0^\infty g^{(1)}\cos\lambda t
=\frac12\int_{\mathbb R}g^{(1)}\cos\lambda t$ by evenness; shifting the contour to
$\operatorname{Im}t=2\pi-\eta$ and letting $\eta\downarrow0$ gives the factor $e^{-2\pi\lambda}$, a double
pole producing the extra factor $(1+2\pi\lambda)$}
\lstep{1}{5}{(Second order.)  $g^{(2)}$ is real-analytic on $[0,\infty)$ with $g^{(2)}(0)=0$ and
$(g^{(2)})'(0)=-\zeta^2/30$.}
\lproof
\lstep{2}{1}{$\dfrac{2\zeta^2S_{t\sigma}^2S_{t(1-\sigma)}^2}{\sinh^2\frac t2
\bigl(\sinh\frac t2+2\zeta S_{t\sigma}S_{t(1-\sigma)}\bigr)}
=\zeta^2\sigma^2(1-\sigma)^2\,t\,\bigl(1+O(t)\bigr)$ as $t\to0$, uniformly in $\sigma\in[0,1]$.}
\lby{$S_x=\frac x2(1+O(x^2))$, so the numerator is
$2\zeta^2\frac{t^4}{16}\sigma^2(1-\sigma)^2(1+O(t^2))$ and the denominator is
$\frac{t^2}4\cdot\frac t2(1+O(t))=\frac{t^3}8(1+O(t))$}
\lstep{2}{2}{$g^{(2)}(t)=-\zeta^2t\int_0^1\sigma^2(1-\sigma)^2\dd\sigma+O(t^2)=-\frac{\zeta^2}{30}t+O(t^2)$.}
\lby{$\langle2\rangle1$, $\int_0^1\sigma^2(1-\sigma)^2\dd\sigma=\frac1{30}$, and dominated convergence}
\lqed{2}{3}\lby{$\langle2\rangle1$, $\langle2\rangle2$ and analyticity from $\langle1\rangle1$}
\lstep{1}{6}{$\displaystyle\operatorname{Re}\int_0^\infty g^{(2)}e^{i\lambda t}\dd t
=-\frac{(g^{(2)})'(0)}{\lambda^2}-\frac1{\lambda^3}\operatorname{Im}\!\int_0^\infty(g^{(2)})'''e^{i\lambda t}\dd t
=\frac{\zeta^2}{30\lambda^2}-\frac1{\lambda^3}\operatorname{Im}\!\int_0^\infty(g^{(2)})'''e^{i\lambda t}\dd t.$}
\lproof
\lstep{2}{1}{Three integrations by parts:
$\int_0^\infty ue^{i\lambda t}\dd t=\frac{iu(0)}\lambda-\frac{u'(0)}{\lambda^2}
-\frac{iu''(0)}{\lambda^3}+\frac1{(i\lambda)^3}\int_0^\infty u'''e^{i\lambda t}\dd t$
for $u\in C^3[0,\infty)$ with $u,u',u''\to0$ at $\infty$ and $u'''\in L^1$.}
\lby{$\int_0^\infty u^{(k)}e^{i\lambda t}=-\frac{u^{(k)}(0)}{i\lambda}-\frac1{i\lambda}
\int_0^\infty u^{(k+1)}e^{i\lambda t}$, iterated three times; the boundary terms at $\infty$ vanish by
$\langle1\rangle1$}
\lstep{2}{2}{With $u=g^{(2)}$ real: $u(0)=0$ and $u''(0)\in\mathbb R$, so
$\operatorname{Re}\bigl[\frac{iu(0)}\lambda-\frac{iu''(0)}{\lambda^3}\bigr]=0$ and
$\operatorname{Re}\frac1{(i\lambda)^3}\int u'''e^{i\lambda t}
=\operatorname{Re}\frac{i}{\lambda^3}\int u'''e^{i\lambda t}
=-\frac1{\lambda^3}\operatorname{Im}\int u'''e^{i\lambda t}$.}
\lby{$\langle1\rangle5$ ($u(0)=0$), $u$ real-valued so $u''(0)$ real, and $(i)^{-3}=i$}
\lqed{2}{3}\lby{$\langle2\rangle1$, $\langle2\rangle2$, $(g^{(2)})'(0)=-\zeta^2/30$}
\lstep{1}{7}{$\bigl|\operatorname{Im}\int_0^\infty(g^{(2)})'''e^{i\lambda t}\dd t\bigr|
\le\|(g^{(2)})'''\|_{L^1(0,\infty)}\le A_3\zeta^2$, and $A_3<\infty$.}
\lby{$|e^{i\lambda t}|=1$; finiteness of $A_3$ from $\langle1\rangle1$ (analyticity at $0$ and
$O(e^{-t/2})$ decay), and the fact that every $t$-derivative of the integrand of \eqref{eq:gsplit}
carries the prefactor $\zeta^2$ and a denominator bounded below by $\sinh^3\frac t2$, so that
$\zeta^{-2}\|(g^{(2)})'''\|_1$ is bounded uniformly for $\zeta\in(0,1]$}
\lstep{1}{8}{Assembling: with $\lambda=\kappa_c/2\pi$,
\[
 \Sigma(\kappa_c)=\frac1{2\pi^2}\Bigl[\frac{\zeta}2\widehat{g^{(1)}}(\lambda)
 +\frac{\zeta^2}{30\lambda^2}-\frac1{\lambda^3}\operatorname{Im}\!\int_0^\infty(g^{(2)})'''e^{i\lambda t}\Bigr],
 \quad\frac1{2\pi^2}\cdot\frac{\zeta^2}{30\lambda^2}=\frac{\zeta^2}{15\kappa_c^2},
\]
and the two remainders are bounded by $\frac1{2\pi^2}\frac{(2\pi)^3}{\kappa_c^3}A_3\zeta^2
=\frac{4\pi A_3\zeta^2}{\kappa_c^3}$ and
$\frac{\zeta}{4\pi^2}A_1(1+\kappa_c)e^{-\kappa_c}$, which is \eqref{eq:taubound}.}
\lby{$\langle1\rangle2$, $\langle1\rangle4$, $\langle1\rangle6$, $\langle1\rangle7$;
$\frac1{2\pi^2}\frac{\zeta^2}{30}\frac{4\pi^2}{\kappa_c^2}=\frac{\zeta^2}{15\kappa_c^2}$}
\lqed{1}{9}\lby{$\langle1\rangle3$ for \eqref{eq:phsym}, $\langle1\rangle8$ for \eqref{eq:taubound};
\eqref{eq:taunu} follows from $\tau=2\tau_+$, $\nu=2\nu_+$, so $5\tau+\nu=12\nu_++10\Sigma$,
and $\frac{10\zeta^2}{15\kappa_c^2}=\frac{2\zeta^2}{3\kappa_c^2}$}

\begin{verdictgap}
The \emph{shape} of \eqref{eq:taubound} --- leading term $\zeta^2/(15\kappa_c^2)$, correction
$O(\zeta^2\kappa_c^{-3})$, first-order contribution exponentially small --- is proved.  The two
constants $A_1,A_3$ are proved finite but their numerical values ($A_3\le0.0374$, printed as $0.036$ before 2026-09-24, see Correction~C10; $A_1\le2.0$) are
obtained by numerical evaluation of an explicit one-dimensional integral (\S\ref{sec:num}), not by a
closed-form estimate; a fully hand-checkable bound would require differentiating \eqref{eq:gsplit}
three times and estimating each of the resulting terms, which we have not done.  This is the only
non-symbolic input to Theorem~\ref{thm:main}.  Note that Theorem~\ref{thm:tau} also \emph{proves}
Note~5, Proposition~4.1 in its integrated form (the pointwise statement $\frac2{15}\zeta^2\kappa^{-3}
+O(\kappa^{-5})$ is not needed anywhere below) and re-derives the exact coefficient
$c_2=-\zeta^2/30$ of \texttt{check\_note5\_theory.tex} independently, from the closed form
\eqref{eq:Fclosed}.
\end{verdictgap}

\section{The tail bound}\label{sec:tail}

\begin{theorem}[ultraviolet tail bound]\label{thm:main}
Let $\kappa_c>0$, $\zeta\in(0,1]$, and put
\[
 \mathcal T:=5\tau+\nu=12\nu_++10\Sigma(\kappa_c),\qquad
 \Phi_\partial:=\frac{\mathcal T}{2(1-\mathcal T)}\quad(\text{assume }\mathcal T<1).
\]
Then $\Phi_W\le\Phi$ and, with $y:=2e^{2\Phi_\partial}\sqrt{\Phi\,\Phi_\partial}$ and $Y:=e^{\Phi_W}y\le e^{\Phi}y$, provided $Y<1$ (corrected 2026-09-24: the hypothesis and the bound were printed with $y$ in place of $Y$, see Correction~C1),
\begin{equation}\label{eq:Psi}
 0\ \le\ \Phi(\zeta)-\Phi_W(\zeta)\ \le\ \Phi_\partial+\frac{Y}{1-Y}
 \qquad\Bigl(\ =\ 2\sqrt{\Phi\,\Phi_\partial}+\Phi_\partial+O\bigl((\Phi+\Phi_\partial)^{3/2}\bigr)\Bigr).
\end{equation}
With Theorem~\ref{thm:tau}, $\mathcal T\le\frac{2\zeta^2}{3\kappa_c^2}+\frac{4.7\zeta^2}{\kappa_c^3}
+\bigl(12+0.51\zeta(1+\kappa_c)\bigr)e^{-\kappa_c}$ (the $4.7$ was printed as $4.6$ before 2026-09-24, see Correction~C10).
\end{theorem}
\lstep{1}{1}{$\Phi_W\le\Phi$.}
\lby{the restriction of $\rho=\omega_Q$ (resp.\ $\sigma=\omega_{\widetilde Q}$) to
$\cF(EH)\subset\cF(I)$ is $\omega_{EQE}$ (resp.\ $\omega_{E\widetilde QE}$), and root fidelity is
nondecreasing under restriction to a subalgebra}
\lstep{1}{2}{$-\log\Fid(\rho'',\sigma)\le\Phi_\partial$.}
\lby{Lemma~\ref{lem:purif} applied to $X=\widetilde Q$, $Y=Q''$, i.e.\ \eqref{eq:purif2}, together
with Theorem~\ref{thm:block} ($\|N\|_2^2\le5\tau+\nu=\mathcal T$) and the monotonicity of
$x\mapsto x/(2(1-x))$; uniform in the finite-dimensional approximation of \S\ref{sec:window},
hence valid in the limit}
\lstep{1}{3}{$e^{-\Phi}=\Fid(\rho,\sigma)\ge\cos\bigl(A(\rho,\rho'')+A(\rho'',\sigma)\bigr)
=e^{-\Phi_W}\Fid(\rho'',\sigma)-\sqrt{1-e^{-2\Phi_W}}\sqrt{1-\Fid(\rho'',\sigma)^2}$,
whenever the sum of angles is $\le\pi/2$.}
\lby{the Bures angle triangle inequality (cited box, \S\ref{sec:triangle}),
$\cos$ decreasing on $[0,\pi/2]$, the addition theorem, and Lemma~\ref{lem:factor}
($\Fid(\rho,\rho'')=e^{-\Phi_W}$)}
\lstep{1}{4}{$\sqrt{1-e^{-2\Phi_W}}\sqrt{1-\Fid(\rho'',\sigma)^2}\le
2e^{\Phi_\partial}\sqrt{\Phi_W\Phi_\partial}\le2e^{\Phi_\partial}\sqrt{\Phi\,\Phi_\partial}$.}
\lby{$1-e^{-2x}\le2x$ for $x\ge0$ applied to $x=\Phi_W$, and
$1-e^{-2z}\le2z\,$ with $z=-\log\Fid(\rho'',\sigma)\le\Phi_\partial$ after writing
$1-\Fid^2=1-e^{-2z}\le2z\le2\Phi_\partial$; the factor $e^{\Phi_\partial}$ absorbs the crude step
$\sqrt{2\Phi_W}\sqrt{2\Phi_\partial}\le2e^{\Phi_\partial}\sqrt{\Phi_W\Phi_\partial}$; then
$\langle1\rangle1$}
\lstep{1}{5}{$e^{-\Phi}\ge e^{-\Phi_W}e^{-\Phi_\partial}-y'=e^{-\Phi_W}\bigl(e^{-\Phi_\partial}-e^{\Phi_W}y'\bigr)$ (corrected 2026-09-24; the printed $e^{-\Phi_W}(e^{-\Phi_\partial}-y')$ does not follow, see Correction~C1) with
$y':=2e^{\Phi_\partial}\sqrt{\Phi\Phi_\partial}$, hence
$\Phi-\Phi_W\le-\log(e^{-\Phi_\partial}-e^{\Phi_W}y')=\Phi_\partial-\log(1-e^{\Phi_W}y'e^{\Phi_\partial})$
and $e^{\Phi_W}y'e^{\Phi_\partial}=e^{\Phi_W}y=Y$; the logarithm is defined because $Y<1$.}
\lby{$\langle1\rangle3$, $\langle1\rangle4$ and $\Fid(\rho'',\sigma)\ge e^{-\Phi_\partial}$}
\lqed{1}{6}\lby{$-\log(1-u)\le u/(1-u)$ with $u=Y<1$, which gives \eqref{eq:Psi};
$\langle1\rangle1$ gives the left inequality}

\begin{corollary}[explicit decay in $\kappa_c$]\label{cor:rate}
Suppose $\Phi(\zeta)\le f\zeta^2$ and $\kappa_c\ge\kappa_\dagger(\zeta):=2\log\frac1\zeta+2\log(1+\kappa_c)+3$
(so that $12e^{-\kappa_c}\le0.6\,\zeta^2\kappa_c^{-2}$) and $\kappa_c\ge10$.  Then
\[
 \mathcal T\le\frac{1.75\,\zeta^2}{\kappa_c^{2}},\qquad
 \Phi-\Phi_W\ \le\ \frac{2\sqrt{0.875\,f}\;\zeta^{2}}{\kappa_c}\,\bigl(1+O(\zeta^2)\bigr)
 \ +\ \frac{0.875\,\zeta^2}{\kappa_c^2}\bigl(1+O(\zeta^2)\bigr).
\]
For the free fermion one may take $f=0.0085$, a numerical input (Note~5, Table~4, printed as ``Table~3'' before 2026-09-24, see Correction~C3; all tabulated $\Phi/\zeta^2\le0.00838$,
and $\lim_{\zeta\to0}\Phi/\zeta^2=1/(12\pi^2)=0.008444$ by \texttt{uhlmann\_upper\_bound.tex} +
\texttt{quadratic\_limit.tex}).  With this $f$ and the numerical values $A_1\le2.0$, $A_3\le0.0374$ of Theorem~\ref{thm:tau} (double precision, not certified; $A_3\le0.036$ before 2026-09-24, see Correction~C10), a numerical evaluation of the corrected chain of Theorem~\ref{thm:main} over the admissible region, not an analytic argument, gives (corrected 2026-09-24, see Correction~C3)
\begin{equation}\label{eq:final}
 \Phi(\zeta)-\Phi_W(\zeta)\ \le\ \frac{0.173\,\zeta^2}{\kappa_c}\;+\;\frac{0.88\,\zeta^2}{\kappa_c^2}.
\end{equation}
\end{corollary}
\lstep{1}{1}{$\mathcal T\le\frac{\zeta^2}{\kappa_c^2}\bigl(\frac23+\frac{4.7}{\kappa_c}\bigr)+\frac{12e^{-3}\zeta^2}{(1+\kappa_c)^2}
+0.51e^{-3/2}\zeta^2e^{-\kappa_c/2}=\frac{\zeta^2}{\kappa_c^2}\,s(\kappa_c)\le\frac{1.75\zeta^2}{\kappa_c^2}$ for $\kappa_c\ge10$ (step corrected 2026-09-24, see Corrections~C2 and~C10).}
\lby{Theorem~\ref{thm:tau} \eqref{eq:taunu}; the hypothesis on $\kappa_\dagger$ gives
$e^{-\kappa_c}\le e^{-3}\zeta^2(1+\kappa_c)^{-2}$, so $12e^{-\kappa_c}\le12e^{-3}\zeta^2(1+\kappa_c)^{-2}\ (\le0.6\zeta^2\kappa_c^{-2})$ and, using its square root $e^{-\kappa_c/2}\le e^{-3/2}\zeta(1+\kappa_c)^{-1}$ once,
$0.51\zeta(1+\kappa_c)e^{-\kappa_c}\le0.51e^{-3/2}\zeta^2e^{-\kappa_c/2}\le0.077\,\zeta^2\kappa_c^{-2}$ for $\kappa_c\ge10$ (the printed ``$\le0.51e^{-3}\zeta^3(1+\kappa_c)^{-1}\le0.003\zeta^2\kappa_c^{-2}$'' is false for moderate $\zeta$, see Correction~C2);
$s(\kappa_c):=\frac23+\frac{4.7}{\kappa_c}+\frac{12e^{-3}\kappa_c^2}{(1+\kappa_c)^2}+0.51e^{-3/2}\kappa_c^2e^{-\kappa_c/2}$ decreases on $[10,\infty)$ (the third term grows with derivative $24e^{-3}\kappa_c(1+\kappa_c)^{-3}<4.7\kappa_c^{-2}$, i.e.\ more slowly than the second term falls, and $\kappa_c^2e^{-\kappa_c/2}$ decreases for $\kappa_c>4$), and $s(10)=\frac23+0.47+0.494+0.077=1.707\le1.75$ (the printed sum $\frac23+0.46+0.6+0.003$ rested on the false $0.003$)}
\lstep{1}{2}{$\Phi_\partial\le\frac{\mathcal T}{2(1-\mathcal T)}\le\frac{0.875\zeta^2}{\kappa_c^2}(1+O(\zeta^2))$.}
\lby{$\langle1\rangle1$ and $\mathcal T\le1.75\zeta^2/100\le0.0175$; hence $\mathcal T<1$, $\Phi_\partial\le0.00891$ and, as $\Phi_W\le\Phi\le f\zeta^2\le f$, the number $Y=e^{\Phi_W}y$ of Theorem~\ref{thm:main} obeys $Y\le e^{f}\,2e^{0.018}\sqrt{0.00891f}\le0.193\,e^{f}f^{1/2}$, which is $<1$ for $f\le1$ and $\le0.018$ for $f=0.0085$, so the hypotheses of Theorem~\ref{thm:main} hold (check added 2026-09-24, see Correction~C3)}
\lqed{1}{3}\lby{Theorem~\ref{thm:main} as corrected (Correction~C1), $\langle1\rangle2$, $\Phi_W\le\Phi\le f\zeta^2$ (the factor $e^{\Phi_W}\le e^{f\zeta^2}=1+O(\zeta^2)$ in $Y$ goes into the factors $1+O(\zeta^2)$), and
$2\sqrt{\Phi\Phi_\partial}\le2\sqrt{f\zeta^2\cdot0.875\zeta^2\kappa_c^{-2}}
=2\sqrt{0.875f}\,\zeta^2/\kappa_c$; with $f=0.0085$, $2\sqrt{0.875\cdot0.0085}=0.1725$.  This proves the display, given the numerical values $A_1\le2.0$, $A_3\le0.0374$ that enter through \eqref{eq:taunu}.  \eqref{eq:final} is \emph{not} obtained from it by dropping the factors $1+O(\zeta^2)$: with $\mathcal T\le1.75\zeta^2/\kappa_c^2$ alone the chain exceeds the right side of \eqref{eq:final} (already the printed chain is $1.032$ times it at $(\zeta,\kappa_c)=(1,10)$).  \eqref{eq:final} rests instead on a numerical evaluation, in double precision, of the corrected chain with the $\mathcal T$ of \eqref{eq:taunu} over the admissible region (edge, corner and grid scans): its ratio to the right side of \eqref{eq:final} has supremum $0.98405$, at the corner $(\zeta,\kappa_c)=(11e^{-3.5},10)$ (corrected 2026-09-24, see Corrections~C3 and~C10; scripts in \texttt{numerics/kb-checks-2026-09-24/} and \texttt{numerics/kb-checks-2026-09-23/})}

\begin{verdictgap}
Proved: a tail bound $\Phi-\Phi_W\le C\zeta^2/\kappa_c+C'\zeta^2/\kappa_c^2$ with finite $C,C'$ whenever $\Phi\le f\zeta^2$; the explicit constants
$C=0.173$, $C'=0.88$ of \eqref{eq:final} rest on numerical inputs ($A_1\le2.0$ and $A_3\le0.0374$, printed as $0.036$ before 2026-09-24 (Correction~C10), in double precision, not certified; $f=0.0085$; a numerical evaluation of the corrected chain), see Corrections~C3 and~C4.  Both are valid once the window reaches
$\kappa_c\ge\max\{10,\ 2\log(1/\zeta)+2\log(1+\kappa_c)+3\}$, the region of Corollary~\ref{cor:rate} (printed here as $\kappa_c\ge2\log(1/\zeta)+2\log\kappa_c+3$ before 2026-09-24).  \emph{Not} proved: the conjectured
$\Phi-\Phi_W\le C\zeta^2\kappa_c^{-2}$ of \texttt{open\_problems\_plan.tex}, \S2.1.  Two losses
separate them.  (i) The Bures triangle inequality contributes the cross term
$2\sqrt{\Phi\Phi_\partial}\sim\zeta^2/\kappa_c$, which is intrinsic to any metric argument: the true
$\Phi-\Phi_W$ is quadratic in the discarded perturbation, a metric triangle inequality is linear in
it.  (ii) The constant $5$ of Theorem~\ref{thm:block} inflates $\Phi_\partial$ by about an order of
magnitude over its true value $\approx\Sigma\approx\zeta^2/(15\kappa_c^2)$ (\S\ref{sec:num}).
Removing (i) would need a genuinely second-order (Fisher-metric) argument comparing $\rho''$ and
$\sigma$ \emph{relative to} $\rho$, e.g.\ a two-sided Bogoliubov extension in which the window and its
complement are treated by a single global purification; we did not find one.
\end{verdictgap}
\section{The box cutoff}\label{sec:box}

The numerics of Note~5, \S6 uses \emph{two} compressions: first onto the box modes
$\mathfrak h_N\subset L^2(I)$ ($\Lambda$, $\kappa_{\max}$; Note~5, \S3.1), then onto the spectral
window $\{\eps<Q_N<1-\eps\}$ of the \emph{compressed} vacuum symbol $Q_N$.  Two consequences.

\paragraph{(a) All tabulated values are lower bounds.}  Both compressions are compressions onto
finite-dimensional subspaces of $H$, so for the resulting projection $P_{\rm sub}$ of rank
$n_{\rm sub}$, Note~5, Theorem~3.1 gives $\Phi_{P_{\rm sub}}\le\Phi$ unconditionally.  This is
independent of everything in this note and is the only part of the error bar that is a theorem
without qualifications (up to the arithmetic of \S\ref{sec:bars}).

\paragraph{(b) The upper bound is \emph{not} directly the one of Theorem~\ref{thm:main}.}
Theorem~\ref{thm:main} bounds $\Phi-\Phi_W$ where $\Phi_W$ is the compression to the spectral window
$E=\one_{[-\kappa_c,\kappa_c]}(K)$ of the \emph{exact continuum} $K$.  The computed quantity is the
compression to $P_{\rm sub}\subset\mathfrak h_N$, and $P_{\rm sub}\ne E$: the box replaces the
continuous modular spectrum by $N$ discrete modes with $|\kappa_j|\le\kappa_{\max}$ and distorts the
occupations near the box edge by $\approx(\pi\Lambda k_j)^{-1}$ (Note~5, \S3.1).  We have \emph{no}
bound on $\Phi_E-\Phi_{P_{\rm sub}}$.  Empirically it is small: Note~5, Table~2 (printed as ``Table~1'' before 2026-09-24, see Correction~C6) varies
$(\Lambda,\kappa_{\max})$ over nine combinations at $\zeta=1/15$ with a spread of $\pm0.06\,\%$ in
$\Phi_{\rm sub}$, and $\Lambda=120$ changes the $\kappa_c=18.4$ value by $-3\times10^{-5}$ relative
(\texttt{numerics/README.md}).  Everything below that uses the upper half of the bar is therefore
labelled ``modulo the box discretisation''; it also rests on the numerical inputs of \eqref{eq:final} and is not certified (printed as ``certified modulo the box discretisation'' before 2026-09-24, see Correction~C8).

\paragraph{(c) When the window is the whole box.}
At $\eps=10^{-14}$ ($\kappa_c=32.2$) and $\Lambda=60$, $\kappa_{\max}=30$ one has
$n_{\rm sub}=N=91$: the spectral window discards \emph{nothing}, and the residual truncation is the
box itself, whose largest modular energy is $\kappa_{\max}^{\rm eff}=29.6$ (the largest $|\kappa_j|$
actually realised).  In this regime the correct second-order tail is not the window one
$\zeta^2/(15\kappa_c^2)=0.0667\,\zeta^2\kappa_c^{-2}$ but the \emph{box} one obtained from the exact
cutoff asymptotics of the Bures--Fisher form,
\begin{equation}\label{eq:boxtail}
 g_B-g_B^{(\Lambda)}=0.513\,r\,\Lambda^{-2}\quad(\texttt{referee\_quadratic\_limit.tex}, \S7)
 \ \Longrightarrow\
 f_2-f_2^{(\kappa_{\max})}=\frac{0.513}{8\,\kappa_{\max}^2}=\frac{0.0641}{\kappa_{\max}^2},
\end{equation}
i.e.\ $0.0641\,\zeta^2\kappa_{\max}^{-2}$, $4\,\%$ below the window value.  At
$\kappa_{\max}^{\rm eff}=29.6$, $\zeta=1/15$ this is $3.25\times10^{-7}$ (against
$3.61\times10^{-7}$ for the window formula at $\kappa_c=32.2$), and
$\Phi_{\rm sub}+{\rm tail}=3.48658\times10^{-5}+3.25\times10^{-7}=3.5191\times10^{-5}$,
consistent with the $\eps=10^{-8}$ and $10^{-11}$ corrected values ($3.5193$, $3.5199$) to
$2\times10^{-4}$ relative.  The ``$+$tail'' column printed in
\texttt{numerics/results\_hp14.txt} uses $\kappa_c$ and is therefore $1.2\times10^{-3}$ relative too
small; the corrected recipe is $\kappa_{\rm eff}=\min(\kappa_c,\kappa_{\max}^{\rm eff})$ with the
coefficient $1/15$ for a window-limited and $0.0641$ for a box-limited run.

\section{Error bars (not certified, see Correction~C8)}\label{sec:bars}

\paragraph{Arithmetic.}  The evaluation of Note~5, Theorem~2.1 inside the window is done in
$40$-digit arithmetic (\texttt{fidelity\_flint.py}), so the arithmetic error is negligible; the
inputs are not exact.  Write $Q_N^{\rm c}$, $\widehat D^{\rm c}$ for the computed matrices and
$\eta:=\max\{\|Q_N-Q_N^{\rm c}\|_1,\|\widehat D-\widehat D^{\rm c}\|_1\}$.  By Lemma~\ref{lem:purif}
and Powers--St\o rmer ($\|A^{1/2}-B^{1/2}\|_2^2\le\|A-B\|_1$ for $A,B\ge0$),
$\|N(X,X')\|_2^2\le\|X-X'\|_2^2+4\|X-X'\|_1\le5\|X-X'\|_1$, so the Bures angle moves by at most
$\arcsin\sqrt{5\eta}$ under either perturbation and
\begin{equation}\label{eq:arith}
 \bigl|\Phi_{\rm sub}-\Phi_{\rm sub}^{\rm c}\bigr|
 \ \le\ \sqrt{2\Phi}\cdot2\sqrt{5\eta}\,(1+o(1))+5\eta .
\end{equation}
With the quoted quadrature accuracies ($3\times10^{-11}$ relative on $\widehat D$, Note~5 \S6;
$\Lambda$-trapezoid with $4\times10^5$ nodes for $Q_N$) a conservative $\eta=10^{-9}$ gives
$|\Phi_{\rm sub}-\Phi^{\rm c}_{\rm sub}|\le2\sqrt{2\cdot3.5\times10^{-5}}\sqrt{5\times10^{-9}}
=1.2\times10^{-6}$ at $\zeta=1/15$ --- i.e.\ $3\,\%$, and it is again the cross term that dominates.
An honest quadrature audit (interval arithmetic on $\widehat D$ and $Q_N$) would reduce $\eta$ to
$\sim10^{-13}$ and this term to $0.4\,\%$; that audit has \emph{not} been carried out here, and this
is the second-largest gap in the certification.

\paragraph{The bar.}  Combining (a), (b), \eqref{eq:final} (whose constants rest on numerical inputs, see Correction~C4, so the upper half below does too) and \eqref{eq:arith}: with
$\kappa_{\rm eff}=\min(\kappa_c,\kappa_{\max}^{\rm eff})$,
\begin{equation}\label{eq:bar}
 \Phi^{\rm c}_{\rm sub}-\eta_{\rm num}\ \le\ \Phi(\zeta)\ \le\
 \Phi^{\rm c}_{\rm sub}+\eta_{\rm num}+\frac{0.173\,\zeta^2}{\kappa_{\rm eff}}
 +\frac{0.88\,\zeta^2}{\kappa_{\rm eff}^2}\quad(\text{upper half: modulo the box, \S\ref{sec:box}b}).
\end{equation}
Table~\ref{tab:bars} lists the resulting intervals together with the \emph{extrapolated} best value
$\Phi^{\rm c}_{\rm sub}+\Sigma$, $\Sigma$ the proved leading tail of Theorem~\ref{thm:tau}
(coefficient $1/15$ or $0.0641$ according to \S\ref{sec:box}c), whose own uncertainty (proved in form; $C_2$ is numerically evaluated, see Correction~C4) is
$C_2\zeta^2\kappa_{\rm eff}^{-3}\le0.47\,\zeta^2\kappa_{\rm eff}^{-3}$ ($\approx24\,\%$ of $\Sigma$: $15\cdot0.47/\kappa_{\rm eff}=0.24$ for the coefficient $1/15$ and $0.47/(0.0641\,\kappa_{\rm eff})=0.25$ for $0.0641$; printed as $0.6\,\%$, with $C_2\le0.46$, before 2026-09-24, see Corrections~C5 and~C10; at
$\kappa_{\rm eff}=29.6$) --- but which is an estimate of $\Phi-\Phi_W$ only to the extent that the
cross term of Theorem~\ref{thm:main} is not saturated.

\begin{table}[ht]\centering\small
\caption{The bar \eqref{eq:bar} (units $10^{-5}$; not certified, since its upper end rests on the numerical inputs of \eqref{eq:final}; printed as ``Certified bar'' before 2026-09-24, see Correction~C8) from the raw sub-compressed values of
\texttt{results\_A2.txt} ($\kappa_c=18.4$), \texttt{results\_hp2.txt} ($25.3$) and
\texttt{results\_hp14.txt} ($32.2$); $\Phi^{\rm c}_{\rm sub}$ is the largest of the three (all are
lower bounds).  $\eta_{\rm num}$ is omitted from the printed interval (add $\pm3\,\%$).
``best'' $=\Phi^{\rm c}_{\rm sub}+\Sigma$ is extrapolated, not certified.}\label{tab:bars}
\begin{tabular}{r r c r r}\toprule
$\zeta$ & $10^5\Phi^{\rm c}_{\rm sub}$ & interval & $10^5\,$best & width/$\Phi$\\\midrule
$1/120$ & 0.05765 & $[0.0576,\ 0.1052]$ & 0.05816 & $82\%$\\
$1/60$  & 0.22871 & $[0.2287,\ 0.4190]$ & 0.23074 & $83\%$\\
$1/30$  & 0.90005 & $[0.9000,\ 1.6609]$ & 0.90817 & $85\%$\\
$1/15$  & 3.48658 & $[3.4866,\ 6.5309]$ & 3.51910 & $87\%$\\
$2/15$  & 13.1093 & $[13.109,\ 25.285]$ & 13.2393 & $93\%$\\
$4/15$  & 46.5158 & $[46.516,\ 104.92]$ & 47.2564 & $126\%$\\
$8/15$  & 151.488 & $[151.49,\ 385.09]$ & 154.450 & $154\%$\\
\bottomrule\end{tabular}\end{table}

\begin{verdictgap}
The bar \eqref{eq:bar} (not certified, see Correction~C8) is a factor $1.8$--$2.5$ wide, i.e.\ far weaker than the heuristic $0.3$--$1\,\%$ of
Note~5 and than the $2\times10^{-4}$ mutual consistency of the three windows.  Two things must improve
before the tabulated digits can be called certified: the cross term of Theorem~\ref{thm:main}
(\S\ref{sec:tail}, verdict) and an interval-arithmetic audit of $\widehat D$ and $Q_N$.  What
\emph{is} now proved and was not before: (i) the tabulated values are lower bounds with a quantified
arithmetic margin; (ii) the correction added to them, $\zeta^2/(15\kappa_c^2)$ (resp.\
$0.0641\zeta^2\kappa_{\max}^{-2}$), is the exact leading ultraviolet occupation tail, with a proved (given the numerically evaluated $C_2$, see Correction~C4)
relative error $\le25\,\%$ at $\kappa_{\rm eff}\simeq30$ ($15\cdot0.47/29.6=0.238$ for the coefficient $1/15$, $0.47/(0.0641\cdot29.6)=0.248$ for $0.0641$; printed as $0.6\,\%$ before 2026-09-24, see Corrections~C5 and~C10), while the measured one is $2.2$--$6.3\,\%$ at $\kappa_c=18.4$ and $0.7$--$1.8\,\%$ at $\kappa_c=32.2$ (\S\ref{sec:num}(4)) --- Note~5's heuristic ``12--56\,\% off''
estimate of the correction was too pessimistic for the \emph{occupation}, and the true uncertainty of
the corrected value lies in the step from the occupation to the fidelity, not in the occupation.
\end{verdictgap}
\section{The rate}\label{sec:rate}

The rate $\Phi=f_2\zeta^2[1+O(\log^{-2}(1/\zeta))]$ claimed by an earlier version of Note~5, Cor.~4.3 (printed as ``Cor.~4.2'' before 2026-09-24; the current version says that this rate was wrong, see Correction~C9), was to follow by combining
(i) the tail bound, (ii) the cutoff error of the Bures--Fisher form
(\texttt{quadratic\_limit.tex}, Prop.~7.1: $0\le g_B-g_B^{(\kappa_c)}\le\frac{2r}{3\kappa_c^2}(1+o(1))$,
exactly $0.513\,r\kappa_c^{-2}$), and (iii) a finite-dimensional cubic estimate
(\texttt{quadratic\_limit.tex}, Problem~7.2):
\begin{equation}\label{eq:cubic}
 \Bigl|\Phi_P(\zeta)-\tfrac{\zeta^2}8g_B^{(P)}\Bigr|\le C\,e^{c\,\kappa_c}\,\zeta^3
 \qquad(P\le E=\one_{[-\kappa_c,\kappa_c]}(K)),
\end{equation}
with $c=1$ the value suggested there by $\min\operatorname{spec}(Q_P)^{-1}\le2e^{\kappa_c}$.

\begin{proposition}[what the tail bound does and does not give]\label{prop:rate}
Assume \eqref{eq:cubic} for some $c>0$.
\begin{enumerate}[label=(\alph*),leftmargin=2em]
\item If $c<\tfrac12$, then choosing $\kappa_c=\theta\log\frac1\zeta$ with any
 $2<\theta<\tfrac1c$ gives
 \[
  \Phi(\zeta)=f_2\,\zeta^2\Bigl[1+O\bigl(\log^{-1}(1/\zeta)\bigr)\Bigr],\qquad f_2=\frac{g_B}8=\frac1{12\pi^2}.
 \]
\item If $c\ge\tfrac12$ --- in particular for the natural value $c=1$ --- the two constraints on
 $\kappa_c$ are incompatible and \emph{no} rate follows from Theorem~\ref{thm:main}
 together with \eqref{eq:cubic}.
\end{enumerate}
\end{proposition}
\lstep{1}{1}{Theorem~\ref{thm:main} forces $\kappa_c>2\log\frac1\zeta$ asymptotically: the term
$\nu\le2e^{-\kappa_c}$ of $\mathcal T$ contributes $\Phi_\partial\ge\tfrac12\nu_+$ and hence, through the
cross term, at least $2\sqrt{\Phi\cdot\tfrac12\nu_+}\approx\sqrt{2f}\,\zeta e^{-\kappa_c/2}$ to the
bound, which is $o(\zeta^2)$ only if $e^{-\kappa_c/2}=o(\zeta)$.}
\lby{Theorem~\ref{thm:main} \eqref{eq:Psi} and Theorem~\ref{thm:tau} \eqref{eq:taunu};
$\nu_+=\log(1+e^{-\kappa_c})$}
\lstep{1}{2}{\eqref{eq:cubic} forces $\kappa_c<\frac1c\log\frac1\zeta$ asymptotically:
$Ce^{c\kappa_c}\zeta^3=o(\zeta^2)$ iff $c\kappa_c-\log\frac1\zeta\to-\infty$.}
\lby{elementary}
\lstep{1}{3}{\lcase $c<\frac12$: pick $\theta\in(2,1/c)$, $\kappa_c=\theta\log\frac1\zeta$.  Then
$e^{-\kappa_c}=\zeta^\theta=o(\zeta^2)$, the hypothesis $\kappa_c\ge\kappa_\dagger(\zeta)$ of
Corollary~\ref{cor:rate} holds for small $\zeta$, and
\[
 \Phi_W\le\Phi\le\Phi_W+\frac{0.173\zeta^2}{\kappa_c}+\frac{0.88\zeta^2}{\kappa_c^2},
 \qquad
 \Bigl|\Phi_W-\frac{\zeta^2}8g_B^{(\kappa_c)}\Bigr|\le C\zeta^{3-c\theta},
 \qquad
 0\le g_B-g_B^{(\kappa_c)}\le\frac{0.52\,r}{\kappa_c^2},
\]
whence $|\Phi-f_2\zeta^2|\le\frac{0.173\zeta^2}{\theta\log(1/\zeta)}+O\bigl(\zeta^2\log^{-2}(1/\zeta)\bigr)
+C\zeta^{3-c\theta}$ and $3-c\theta>2$.}
\lby{Theorem~\ref{thm:main}, Corollary~\ref{cor:rate} (only the finiteness of its constants is needed; the values $0.173$, $0.88$ rest on numerical inputs, see Correction~C4), \eqref{eq:cubic} applied to $P\uparrow E$
(the supremum over $P\le E$ of $\Phi_P$ is $\Phi_W$, and $\sup_Pg_B^{(P)}=g_B^{(\kappa_c)}$ by
definition), and \texttt{quadratic\_limit.tex}, Prop.~7.1}
\lqed{1}{4}\lby{$\langle1\rangle1$, $\langle1\rangle2$: the intervals
$(2\log\frac1\zeta,\infty)$ and $(0,\frac1c\log\frac1\zeta)$ are disjoint for $c\ge\frac12$ and
$\zeta$ small; $\langle1\rangle3$ for (a)}

\begin{verdictgap}
The claimed $\log^{-2}$ rate is \emph{not} obtained; the best this route can give is $\log^{-1}$, and
only under a strengthening of Problem~7.2 to an exponent $c<1/2$.  The obstruction is the same cross
term as in \S\ref{sec:tail}: it converts the requirement ``the window must reach past
$\kappa_*(\zeta)\approx2\log(1/\zeta)$, where the vacuum occupation drops below the recovered one''
(which is genuine, cf.\ Note~5, Cor.~4.3, printed as ``Cor.~4.2'' before 2026-09-24, see Correction~C9) into a \emph{linear} rather than quadratic penalty.  A
version of Theorem~\ref{thm:block} in which the vacuum tail $\nu$ enters through the Hellinger
combination $\sum_{|\kappa|>\kappa_c}(\sqrt{\widetilde q_\kappa}-\sqrt{q_\kappa})^2$ instead of
$\sum(q_\kappa+\widetilde q_\kappa)$ would remove the constraint $\kappa_c>2\log(1/\zeta)$ altogether;
on the diagonal this is elementary, and the whole difficulty is the off-diagonal block $X_{12}$, for
which the square root is not Lipschitz.  Note also that Note~5's own heuristic for the $\log^{-2}$
rate estimates $\Phi-\Phi_{\kappa_*}$, not $\Phi-\Phi_{\kappa_c}$ for a free $\kappa_c$, so it is not
in contradiction with (b).
\end{verdictgap}
\section{Numerical verification}\label{sec:num}

Scripts: \texttt{numerics/pt\_checks.py} (the two constants and $\Sigma$),
\texttt{numerics/pt\_block.py} (Theorem~\ref{thm:block} on the free-fermion data),
\texttt{numerics/pt\_lemmas.py}, \texttt{numerics/pt\_lem3.py} (Lemmas~\ref{lem:factor},
\ref{lem:purif}), \texttt{numerics/pt\_bars.py} (Table~\ref{tab:bars}),
\texttt{rigor/pt\_tailconst.py} ($A_3$).  Interpreter:
the project venv (numpy/mpmath/python-flint).

\paragraph{(1) Lemma~\ref{lem:purif} (purification bound).}  $40$ random non-commuting pairs
$0.02\le X,Y\le0.98$ on $\mathbb C^3$, $40$-digit arithmetic: $\max\Fid_{\rm purif}/\Fid_{\rm exact}
=1.000$ (never exceeded; $\Fid_{\rm exact}$ from Note~5, Theorem~2.1), typical ratio $0.70$--$0.88$.
$30$ perturbative pairs ($\|Y-X\|\sim10^{-3}$) on $\mathbb C^4$:
$\max\Phi_{\rm exact}/\Phi_{\rm purif}=0.960$, i.e.\ the bound is tight to $4\,\%$ exactly in the
regime in which it is used.  \emph{(A double-precision version of the same test reports spurious
violations once $\min\operatorname{spec}\lesssim10^{-9}$; the check must be done in extended
precision.)}

\paragraph{(2) Lemma~\ref{lem:factor} (factorisation).}  Random $8\times8$ diagonal $Q$ and
$\widetilde Q=Q-\delta$ with a $5{+}3$ block split:
$-\log\Fid(\rho,\rho'')$ and $\Phi_W$ agree to $5\times10^{-9}$ relative (double precision).

\paragraph{(3) The constants of Theorem~\ref{thm:tau}.}  $\widehat{g^{(1)}}(\kappa_c/2\pi)$ divided by
$(1+\kappa_c)e^{-\kappa_c}$ equals $1.9954,\,1.9994,\,2.0000,\,2.0000,\,1.9997$ at
$\kappa_c=6,8,12,18.4,25.3$ (beyond that the double-precision floor is reached): numerically $A_1=2$ (inferred from these samples, none below $\kappa_c=6$; not certified, see Correction~C4), so
$C_1=1/(2\pi^2)=0.05066$.  This reproduces the independent quadrature of
\texttt{check\_note5\_theory.tex}, \S6.4, which finds
$\widetilde Q^{(1)}(\kappa,\kappa)-q_\kappa=0.00337737\,\kappa e^{-\kappa}$ at $\zeta=1/15$:
$0.00337737\times15=0.050661=1/(2\pi^2)$.  For $A_3$: $\zeta^{-2}\|(g^{(2)})'''\|_{L^1}
=0.0355,\,0.0326,\,0.0275,\,0.0214$ at $\zeta=0.05,0.2,0.5,1$ on the grid of \texttt{rigor/pt\_tailconst.py}; they decrease in $\zeta$, so the supremum $A_3$ is the limit $\zeta\to0$, which was not sampled: $0.0364$ on the same grid and $0.0372$ on a refined grid (printed as ``(monotone decreasing in $\zeta$), so'' before 2026-09-24, see Correction~C10), so
$A_3\le0.0374$ and $C_2=4\pi A_3\le0.47$ (printed as $A_3\le0.036$ and $C_2\le0.46$ before 2026-09-24).

\paragraph{(4) The tail $\Sigma(\kappa_c)$.}  Exact evaluation of
$\Sigma=\frac1{2\pi^2}\int_0^\infty g\cos(\lambda t)\dd t$ (period-by-period Gauss--Legendre):

\begin{center}\small\begin{tabular}{r rrr rrr}\toprule
& \multicolumn{3}{c}{$\Sigma(\kappa_c)/\bigl[\zeta^2/(15\kappa_c^2)\bigr]$}
& \multicolumn{3}{c}{$\bigl[\Sigma-\zeta^2/(15\kappa_c^2)\bigr]\kappa_c^3/\zeta^2$ (bound: $\le0.47$, Correction~C10)}\\
$\zeta$ & $\kappa_c{=}18.4$ & $25.3$ & $32.2$ & $18.4$ & $25.3$ & $32.2$\\\midrule
$1/15$ & 1.0631 & 1.0306 & 1.0184 & 0.0774 & 0.0515 & 0.0394\\
$0.2$  & 1.0608 & 1.0298 & 1.0179 & 0.0745 & 0.0503 & 0.0385\\
$1.0$  & 1.0224 & 1.0114 & 1.0069 & 0.0274 & 0.0192 & 0.0149\\\bottomrule
\end{tabular}\end{center}
The remainder bound $0.47\,\zeta^2\kappa_c^{-3}$ (proved in form; the constant is the numerically evaluated $C_2$, see Corrections~C4 and~C10) holds with a factor $6$--$31$ to spare, and
the measured remainder decays faster than $\kappa_c^{-3}$ (consistent with the next term being
$O(\kappa_c^{-4})$).

\paragraph{(5) Theorem~\ref{thm:block} on the free-fermion data.}  $\widehat D$ from
\texttt{Dhat\_exact\_s0.2\_*} ($\zeta=1/15$), $Q_N$ from \texttt{compression\_box.py}, everything
transported to the eigenbasis of $Q_N$ (so $E$ is a coordinate projection):

\begin{center}\small\begin{tabular}{ll r rrr rr}\toprule
$(\Lambda,\kappa_{\max})$ & $\kappa_c$ & $n_{\rm out}$ & $\tau_+$ & $\tau_-$ & $\nu$
& $\|N\|_2^2$ & $\|N\|_2^2/(5\tau+\nu)$\\\midrule
$(60,30)$  & 18.4 & 30  & 5.743e$-$7 & 5.743e$-$7 & 2.27e$-$8 & 2.215e$-$6 & 0.384\\
$(60,30)$  & 25.3 & 10  & 1.327e$-$7 & 1.327e$-$7 & 2.26e$-$11& 5.279e$-$7 & 0.398\\
$(120,45)$ & 18.4 & 156 & 7.941e$-$7 & 7.498e$-$7 & 6.65e$-$8 & 3.092e$-$6 & 0.397\\
$(120,45)$ & 25.3 & 114 & 3.087e$-$7 & 2.644e$-$7 & 7.35e$-$11& 1.328e$-$6 & 0.463\\
$(120,45)$ & 32.2 & 71  & 1.119e$-$7 & 7.477e$-$8 & 5.95e$-$14& 5.649e$-$7 & 0.605\\\bottomrule
\end{tabular}\end{center}
\eqref{eq:block} holds throughout with a factor $1.7$--$2.6$ to spare.  Two further checks:
(i) $\tau_+=\tau_-$ to all displayed digits at $\Lambda=60$ (the exact particle--hole symmetry
$\langle1\rangle3$ of Theorem~\ref{thm:tau}); the $6\,\%$ asymmetry at $\Lambda=120$ is a box-edge
artefact.  (ii) The box occupation tail is the \emph{difference} of two continuum tails: at
$\kappa_c=18.4$, $\Lambda=60$, $\kappa_{\max}^{\rm eff}=29.6$ one measures
$\tau_+-\nu_+=5.63\times10^{-7}$ against
$\Sigma(18.4)-\Sigma(29.6)=9.31\times10^{-7}-3.45\times10^{-7}=5.86\times10^{-7}$ ($4\,\%$, box-edge),
confirming both Theorem~\ref{thm:tau} and \S\ref{sec:box}c.
Finally $\|X_{12}\|_2^2=1.02,\,0.48,\,0.21$ ($\times10^{-6}$) and $\|E^\perp h(X)E\|_2^2
=5.1,\,1.8,\,0.65$ ($\times10^{-7}$) at $\kappa_c=18.4,25.3,32.2$ ($\Lambda=120$), i.e.\ both decay
like $\kappa_c^{-2.5}$, \emph{not} faster than the diagonal tail --- which is why the alternative
route through the block-diagonal intermediate state $\omega_{X_{11}\oplus X_{22}}$ (for which the
diagonal part of $\Phi-\Phi_W$ carries \emph{no} cross term) does not improve the exponent either.

\paragraph{(6) The assembled bound.}  At $\zeta=1/15$, $\kappa_c=32.2$: \eqref{eq:final} gives
$\Phi-\Phi_W\le2.77\times10^{-5}$, while the measured difference between the $\kappa_c=32.2$ and the
$\kappa_c=18.4$ values is $5.5\times10^{-7}$ and the predicted total remainder is
$\Sigma=3.3\times10^{-7}$: the bound \eqref{eq:final} (whose constants rest on numerical inputs, see Correction~C4; printed as ``the certified bound'' before 2026-09-24) is a factor $\approx80$ above the truth, entirely
because of the cross term ($2\sqrt{\Phi\Phi_\partial}$ with the proved $\Phi_\partial$; using the
\emph{measured} $\|N\|_2^2$ instead would give $9\times10^{-6}$, still a factor $27$).
\section{Summary of findings}\label{sec:findings}

\begin{enumerate}[leftmargin=1.6em]
\item \textbf{Proved in shape; the constants rest on numerical inputs} (printed as ``Proved (unconditional)'' before 2026-09-24, see Correction~C4).  $\Phi-\Phi_W\le0.173\,\zeta^2\kappa_c^{-1}+0.88\,\zeta^2\kappa_c^{-2}$
for $\kappa_c\ge\max\{10,\ 2\log\frac1\zeta+2\log(1+\kappa_c)+3\}$ (Theorem~\ref{thm:main},
Corollary~\ref{cor:rate}); proved is the form $C\zeta^2\kappa_c^{-1}+C'\zeta^2\kappa_c^{-2}$ with finite $C,C'$, while $0.173$ and $0.88$ rest on $A_1\le2.0$, $A_3\le0.0374$ (double precision, not certified), $f=0.0085$ and a numerical evaluation of the corrected chain (Corrections~C1, C3, C4 and~C10).  The route is: an exact factorisation making $\Phi_W$ the Bures distance to
an explicit third state (Lemma~\ref{lem:factor}); a canonical-purification lower bound for the
fidelity (Lemma~\ref{lem:purif}); and a block estimate (Theorem~\ref{thm:block}) which bounds the
\emph{entire} discarded block --- diagonal, off-diagonal and infrared--ultraviolet --- by the
ultraviolet occupation $\tau$, using nothing but $\widetilde Q\ge0$, operator concavity of
$\sqrt{t(1-t)}$ and a Schur complement.
\item \textbf{Proved, with numerically evaluated constants} (the values $0.47$ and $0.051$ below rest on $A_3\le0.0374$ and $A_1\le2.0$ in double precision, not certified; see Corrections~C4 and~C10).  The occupation tail: $\tau_+=\tau_-$ exactly, and
$\bigl|\int_{\kappa_c}^\infty(\widetilde Q(\kappa,\kappa)-q_\kappa)\dd\kappa-\frac{\zeta^2}{15\kappa_c^2}\bigr|
\le0.47\,\zeta^2\kappa_c^{-3}+0.051\,\zeta(1+\kappa_c)e^{-\kappa_c}$ (Theorem~\ref{thm:tau}; $0.46$ before 2026-09-24, Correction~C10).  This
upgrades Note~5, Proposition~4.1 (a sketch with an uncontrolled remainder and a sign fixed a
posteriori) to a theorem in the only form in which it is used, and re-derives $c_2=-\zeta^2/30$ from
the closed form.
\item \textbf{Not proved.}  The conjectured $\Phi-\Phi_W\le C\zeta^2\kappa_c^{-2}$: the Bures triangle
inequality contributes an irreducible cross term $2\sqrt{\Phi\,\Phi_\partial}=\Theta(\zeta^2/\kappa_c)$.
\item \textbf{Not proved, and now known to fail on this route.}  The rate.  With
Problem~7.2 in its stated form ($Ce^{\kappa_c}\zeta^3$) the two requirements
$\kappa_c>2\log\frac1\zeta$ (from the tail bound) and $\kappa_c<\log\frac1\zeta$ (from the cubic
estimate) are incompatible; a strengthening to $Ce^{c\kappa_c}\zeta^3$ with $c<\frac12$ yields
$\Phi=f_2\zeta^2[1+O(\log^{-1}(1/\zeta))]$ --- still one logarithm short of the $\log^{-2}$ rate claimed by an earlier version of Note~5, Cor.~4.3 (printed as ``Cor.~4.2'' before 2026-09-24; the current version says that this rate was wrong, see Correction~C9)
(Proposition~\ref{prop:rate}).
\item \textbf{Numerics} (printed as ``Certified numerics'' before 2026-09-24, see Correction~C8).  Table~\ref{tab:bars}: the tabulated $\Phi_{\rm sub}$ are rigorous
lower bounds (up to an arithmetic margin of $\approx3\,\%$, dominated again by a cross term and
removable by an interval-arithmetic audit of $\widehat D$ and $Q_N$); the upper bound, which is not certified (see Correction~C8), is a
factor $1.8$--$2.5$ above, and is in addition conditional on the numerical inputs of \eqref{eq:final} (see Correction~C4) and on the box discretisation
(\S\ref{sec:box}b).  The \emph{correction} $\zeta^2/(15\kappa_c^2)$ that Note~5 adds is now proved
correct to $24\,\%$ at $\kappa_c\simeq30$, given the numerically evaluated $C_2$ (measured: $0.7$--$1.8\,\%$ at $\kappa_c=32.2$; printed as $0.6\,\%$ before 2026-09-24, see Corrections~C4, C5 and~C10); what is not proved is that the correction to $\Phi$ equals
the correction to the occupation.
\item \textbf{Correction to the numerics.}  For runs in which the spectral window contains all box
modes ($\eps=10^{-14}$, $\Lambda=60$, $\kappa_{\max}=30$) the tail must be computed with
$\kappa_{\rm eff}=\kappa_{\max}^{\rm eff}=29.6$ and the \emph{box} coefficient
$0.0641=0.513/8$ (\S\ref{sec:box}c), not with $\kappa_c=32.2$ and $1/15$; the printed
``$+$tail'' column of \texttt{results\_hp14.txt} is $1.2\times10^{-3}$ relative too small.
\end{enumerate}

\paragraph{Where to attack next.}  The single obstruction to both the $\kappa_c^{-2}$ tail and the
rate is the same: a bound on $\|\sqrt{\widetilde Q}-\sqrt{X_{11}\oplus X_{22}}\|_2$ (equivalently, on
the Bures angle between $\omega_{\widetilde Q}$ and its pinching) that is \emph{quadratic} in the
off-diagonal block $X_{12}$ rather than governed by $\|X_{12}\|_1^{1/2}$.  On the diagonal the
required inequality is the elementary Hellinger one; the square root's failure to be Lipschitz at the
edge of the spectrum is the whole difficulty, and the Schur-complement constraint
$X_{12}^*X_{11}^{-1}X_{12}\le X_{22}$ used in Theorem~\ref{thm:block}$\langle1\rangle6$ is probably the
right tool for it.

\section{Corrections of 2026-09-24}

These corrections implement knowledge-base card \texttt{tail-bound}: its refereed statement and the
DOCUMENT ERROR notes on this document, together with the referee reports REF-DEP-CFT-B2, B5, B6, B9,
B10 and B11 recorded there; items C5--C10 were added later the same day at the request of the
knowledge-base orchestrator, C10 after the referee report REF-DOC-TAIL.  Every changed passage points to one of the items below, except the
title (line~5) and the header row of Table~\ref{tab:bars} (line~631), which are recorded in C8.  No line
above this section was added or removed, so line numbers cited elsewhere still hold, and the line
numbers below refer to this file.  The check scripts are in \nolinkurl{numerics/kb-checks-2026-09-23/}
and, for C10, in \nolinkurl{numerics/kb-checks-2026-09-24/} of this repository.  Each has its stored
printout (\texttt{.out}) and is described in its folder's \texttt{README.md}.  Run them from the
repository root with \texttt{python3}; they need only the standard library, except
\texttt{a3\_refined.py}, which needs mpmath.  The scripts of 2026-09-23 use the old constant
$C_2\le0.46$; C10 gives the values for $C_2\le0.47$.

\paragraph{C1. Theorem~\ref{thm:main}: step $\langle1\rangle5$ and \eqref{eq:Psi} (lines 57, 464,
466, 494, 496, 497, 499).}
\emph{Before:} step $\langle1\rangle5$ read $e^{-\Phi}\ge e^{-\Phi_W}(e^{-\Phi_\partial}-y')$, and
\eqref{eq:Psi} read $\Phi-\Phi_W\le\Phi_\partial+y/(1-y)$ under the hypothesis $y<1$, where
$y=2e^{2\Phi_\partial}\sqrt{\Phi\Phi_\partial}$.
\emph{Now:} $\langle1\rangle3$, $\langle1\rangle4$ and $\Fid(\rho'',\sigma)\ge e^{-\Phi_\partial}$
give only
\[
 e^{-\Phi}\ \ge\ e^{-\Phi_W}e^{-\Phi_\partial}-y'\ =\ e^{-\Phi_W}\bigl(e^{-\Phi_\partial}-e^{\Phi_W}y'\bigr),
\]
hence $\Phi-\Phi_W\le\Phi_\partial-\log(1-Y)\le\Phi_\partial+Y/(1-Y)$ with $Y:=e^{\Phi_W}y\le e^{\Phi}y$,
provided $Y<1$ (and $\mathcal T<1$ as before).  The printed step subtracted $e^{-\Phi_W}y'$ instead
of $y'$.  Since $e^{-\Phi_W}\le1$, that is too little.  The leading-order form in \eqref{eq:Psi}
does not change, because $Y-2\sqrt{\Phi\Phi_\partial}=O((\Phi+\Phi_\partial)^2)$.  The overview item
(A) now states the hypotheses.  Here $Y$ is a number; it is not the operator $Y=Q''$ to which
Lemma~\ref{lem:purif} is applied in $\langle1\rangle2$, nor the $Y$ of \S\ref{sec:delta}.
\emph{Why:} the printed bound does not follow from the ingredients of the proof.  They allow
equality in the angle triangle inequality at $\langle1\rangle3$, with
$\cos A(\rho,\rho'')=e^{-\Phi_W}$ and $\cos A(\rho'',\sigma)=e^{-\Phi_\partial}$.  For
$\Phi_W=0.0085$ and $\Phi_\partial=10^{-8}$ this configuration gives
$\Phi-\Phi_W=1.0032\,\bigl(\Phi_\partial+y/(1-y)\bigr)$, and as $\Phi_\partial\to0$ the ratio tends
to $\sqrt{(e^{2\Phi_W}-1)/(2\Phi_W)}=1.0043$.  The corrected bound holds there, with ratio $0.9947$
(\texttt{refb6\_psi\_check3.py}).  Source: card \texttt{tail-bound}, the DOCUMENT ERROR note on
$\langle1\rangle5$ (from REF-DEP-CFT-B6, confirmed by REF-DEP-CFT-B9).

\paragraph{C2. Corollary~\ref{cor:rate}, step $\langle1\rangle1$ (lines 517, 518, 520--522).}
\emph{Before:} $0.51\zeta(1+\kappa_c)e^{-\kappa_c}\le0.51e^{-3}\zeta^3(1+\kappa_c)^{-1}
\le0.003\,\zeta^2\kappa_c^{-2}$, and $\frac23+0.46+0.6+0.003\le1.75$.
\emph{Why this was wrong:} the first inequality holds, but the second is false for moderate $\zeta$.
On the admissible region the ratio $0.51\zeta(1+\kappa_c)e^{-\kappa_c}/(\zeta^2\kappa_c^{-2})$
reaches $0.0767$ at the corner $(\zeta,\kappa_c)=(11e^{-3.5},10)$, where $\kappa_c=\kappa_\dagger(\zeta)$
($0.0766$ at the grid point $(0.3325,10)$, \texttt{tailcheck\_ext.py}).  With $0.077$ in place of
$0.003$ and the crude $0.6$, the sum would be $1.804>1.75$.
\emph{Now:} the hypothesis $\kappa_c\ge\kappa_\dagger(\zeta)$ gives
$e^{-\kappa_c}\le e^{-3}\zeta^2(1+\kappa_c)^{-2}$ and its square root
$e^{-\kappa_c/2}\le e^{-3/2}\zeta(1+\kappa_c)^{-1}$.  Hence
$12e^{-\kappa_c}\le12e^{-3}\zeta^2(1+\kappa_c)^{-2}$ and
$0.51\zeta(1+\kappa_c)e^{-\kappa_c}\le0.51e^{-3/2}\zeta^2e^{-\kappa_c/2}$, so
$\mathcal T\le\zeta^2\kappa_c^{-2}s(\kappa_c)$ with
\[
 s(\kappa_c)=\frac23+\frac{4.7}{\kappa_c}+\frac{12e^{-3}\kappa_c^2}{(1+\kappa_c)^2}
 +0.51e^{-3/2}\kappa_c^2e^{-\kappa_c/2}
\]
(with $4.7=10\,C_2$ from C10; $4.6$ before).  $s$ decreases on $[10,\infty)$: the third term grows
with derivative $24e^{-3}\kappa_c(1+\kappa_c)^{-3}<4.7\kappa_c^{-2}$, and $\kappa_c^2e^{-\kappa_c/2}$
decreases for $\kappa_c>4$.  Since $s(10)=0.667+0.47+0.494+0.077=1.707\le1.75$, the statement
$\mathcal T\le1.75\zeta^2/\kappa_c^2$ of the corollary is unchanged.  The value $1.707$ is attained
at the corner, where the condition on $\kappa_\dagger$ holds with equality.  Numerically
$\mathcal T\kappa_c^2/\zeta^2=1.7071$ at the corner and at most $1.7060$ on the grid
(\texttt{refb10\_tail\_sup\_c2.py}); with the old $4.6$ these were $1.6971$ and $1.6960$
(\texttt{refb10\_tail\_sup.py}, \texttt{tailcheck\_ext.py}).  Source: card \texttt{tail-bound},
DOCUMENT ERROR note (i) of AUTH-CFT-B, with the arithmetic of REF-DEP-CFT-B5.

\paragraph{C3. Corollary~\ref{cor:rate}: what \eqref{eq:final} rests on (lines 510, 512, 524, 525,
527).}
\emph{Before:} \eqref{eq:final} was obtained from the display of the corollary by dropping its
factors $1+O(\zeta^2)$ and rounding $0.1725$ to $0.173$ and $0.875$ to $0.88$.  The value $f=0.0085$
was attributed to Note~5, Table~3, and the hypothesis $y<1$ of Theorem~\ref{thm:main} was not checked.
\emph{Why this was wrong:} the factors cannot simply be dropped.  If $\Phi_\partial$ is bounded only
through $\mathcal T\le1.75\zeta^2/\kappa_c^2$, the chain of Theorem~\ref{thm:main} exceeds the right
side of \eqref{eq:final}.  Already the printed chain is $1.032$ times it at $(\zeta,\kappa_c)=(1,10)$
(\texttt{tailcheck.py}, \texttt{tailcheck\_ext.py}), and the corrected chain of C1 is larger still.
The table is Table~4 (\texttt{tab:scan}) of the current Note~5.
\emph{Now:} the display with the factors $1+O(\zeta^2)$ is proved as before, given the numerical $A_1$, $A_3$ of \eqref{eq:taunu}; the factor
$e^{\Phi_W}\le e^{f\zeta^2}$ in $Y$ goes into those factors.  Step $\langle1\rangle2$ checks the
hypotheses of Theorem~\ref{thm:main}: $\mathcal T\le0.0175$, and $Y\le0.193\,e^{f}f^{1/2}$, which is
$<1$ for $f\le1$ and at most $0.018$ for $f=0.0085$.

\eqref{eq:final} is now stated as resting on numerical inputs and a numerical evaluation.  The inputs
are $A_1\le2.0$ and $A_3\le0.0374$ (C10; $0.036$ before) of Theorem~\ref{thm:tau}, evaluated in
double precision and not certified (\S\ref{sec:num}(3)), and $f=0.0085$, a numerical input that the
corollary reads off the tabulated values of Note~5, Table~4, and the limit $1/(12\pi^2)$.  Those
values do not bound $\Phi$ between the grid points or near $\zeta=1$.  The basis recorded on
knowledge-base card \texttt{phi-below-0085-zeta2-numerical} is different: $\Phi\le0.0085\,\zeta^2$
is rigorous on $(0,0.8868]$ by \texttt{rate\_of\_remainder.tex}, Proposition \texttt{prop:rrr}, and
numerical on $[0.8868,1]$, through the monotonicity of $\Phi$ (\texttt{check\_note5\_theory.tex},
\texttt{prop:mono}) and the computed value $\Phi(1.0667)=4.41\times10^{-3}$.  For the evaluation, take $(\zeta,\kappa_c)$ in the
admissible region $0<\zeta\le1$, $\kappa_c\ge10$, $\kappa_c\ge\kappa_\dagger(\zeta)$.  Let
$\bar{\mathcal T}$ be the right side of \eqref{eq:taunu}, and put
$\bar\Phi_\partial=\bar{\mathcal T}/(2(1-\bar{\mathcal T}))$,
$\bar y=2e^{2\bar\Phi_\partial}\sqrt{f\zeta^2\bar\Phi_\partial}$ and $\bar Y=e^{f\zeta^2}\bar y$.
The corrected bound of Theorem~\ref{thm:main} increases with $\mathcal T$, $\Phi$ and $\Phi_W$, so
$\Phi-\Phi_W\le B:=\bar\Phi_\partial+\bar Y/(1-\bar Y)$.

The ratio of $B$ to the right side of \eqref{eq:final} was evaluated in double precision with
$C_2\le0.47$ (C10) by \texttt{numerics/kb-checks-2026-09-24/refb10\_tail\_sup\_c2.py}.  It scans the
edge $\kappa_c=10$, the edge $\kappa_c=\kappa_\dagger(\zeta)$ up to $\kappa_c=30$ and a patch at the
corner.  It also scans the grid of \texttt{tailcheck\_ext.py}: $5.25\times10^6$ admissible points
with $10^{-8}\le\zeta\le1$ and $10\le\kappa_c\le200$, whose maximum is $0.98365$ at $(0.3325,10)$.
Over all scanned points the ratio is largest at the corner
$(\zeta,\kappa_c)=(11e^{-3.5},10)=(0.33217,10)$, where $\kappa_c=\kappa_\dagger(\zeta)=10$, with value
$0.98405$.  The printed chain, with $y$ in place of $Y$, gives $0.98343$ there.  With the old
$C_2\le0.46$ the scripts of 2026-09-23 gave the following.  For the corrected chain, $0.98019$ at
the corner (\texttt{refb10\_tail\_sup.py}) and $0.97979$ on the grid
(\texttt{refb6\_tailcheck\_corrected.py}).  For the printed chain, $0.97957$ at the corner
(\texttt{refb10\_tail\_sup.py}) and $0.97917$ on the grid (\texttt{tailcheck.py},
\texttt{tailcheck\_ext.py}).  The card records the corner as the supremum over the whole admissible
region.  A referee scan up to $\kappa_c=1400$ (REF-DEP-CFT-B11), made with $C_2\le0.46$, found the
ratio decreasing along $\kappa_c=\kappa_\dagger(\zeta)$, to $0.848$ at $\kappa_c=1400$.

With $f=0.008516$ the supremum is $0.98467$ ($0.98080$ with $C_2\le0.46$).  This $f$ comes from the bound
$\Phi\le-\frac12\log(1-2f_2\zeta^2)\le0.008516\,\zeta^2$ on $(0,1]$ of \texttt{rate\_of\_remainder.tex},
Proposition \texttt{prop:rrr}, which could replace the numerical input $f$, with
$2\sqrt{0.875\cdot0.008516}=0.1726\le0.173$.  All of this is a numerical verification of
\eqref{eq:final} given the inputs, not an analytic proof.  Source: card \texttt{tail-bound}, the
DOCUMENT ERROR note of AUTH-CFT-B (its item (ii) and its remark on the table number), its Statement
(b), and the referee reports REF-DEP-CFT-B2, B10 and B11.

\paragraph{C4. The status of the explicit constants (lines 66, 69, 530--532, 610, 619, 648, 700,
745, 764, 798, 804, 806, 812, 829, 831).}
\emph{Before:} line~66 called $C_1\le0.051$ and $C_2\le0.46$ ``certified numerically''.  The verdict
after Corollary~\ref{cor:rate} read ``Proved: an unconditional tail bound \ldots\ with the explicit
constant $C=0.173$'', valid for $\kappa_c\ge2\log(1/\zeta)+2\log\kappa_c+3$.  The summary called
\eqref{eq:final} ``Proved (unconditional)'' and called Theorem~\ref{thm:tau} with the constants $0.46$
and $0.051$ ``Proved''.  Line~764 read ``the proved remainder bound $0.46\,\zeta^2\kappa_c^{-3}$'',
and line~798 called \eqref{eq:final} ``the certified bound''.
\emph{Now:} what is proved is the shape: $A_1$ and $A_3$ are finite, and
$\Phi-\Phi_W\le C\zeta^2/\kappa_c+C'\zeta^2/\kappa_c^2$ with finite $C,C'$ whenever $\Phi\le f\zeta^2$.
The values $A_1\le2.0$ and $A_3\le0.0374$, hence $C_1\le0.051$ and $C_2\le0.47$ ($A_3\le0.036$ and
$C_2\le0.46$ before C10), are double-precision and finite-difference evaluations, not certified; the verdict after Theorem~\ref{thm:tau} already says
that they are numerical.  $f=0.0085$ is a numerical input.  So the constants $0.173$ and $0.88$ of
\eqref{eq:final} rest on these inputs and on the evaluation of C3.  The verdict after
Corollary~\ref{cor:rate} now gives the region of the corollary,
$\kappa_c\ge\max\{10,\,2\log(1/\zeta)+2\log(1+\kappa_c)+3\}$; the printed $2\log\kappa_c$, without
$\kappa_c\ge10$, described a larger region than the corollary covers.  The later uses of
\eqref{eq:final} now say that they inherit these inputs: the bar \eqref{eq:bar} (line~610),
\S\ref{sec:num}(6) (line~798) and the summary (line~829).  Lines~619, 648 and~831 now say that the
remainder of Theorem~\ref{thm:tau} is proved in form with the numerically evaluated $C_2$; the
percentages quoted there are corrected in C5.  Overview item (C) (line~69) said ``with an explicit
$C$''; it now says that $C$ is finite and that the explicit value rests on numerical inputs.  The
proof of Proposition~\ref{prop:rate} (line~700) now says that it needs only the finiteness of the
constants of Corollary~\ref{cor:rate}.  Line~745 said ``$A_1=2$ exactly''; it now says that
$A_1=2$ is numerical, inferred from samples at $\kappa_c\ge6$, and not certified.
Source: card \texttt{tail-bound}, its Statement (b) and its title (``explicit constants 0.173, 0.88
rest on numerical inputs''), its dependency finding and the referee report REF-DEP-CFT-B2, and card
\texttt{tail-constants-numerical}.

\paragraph{C5. The size of the $C_2$ remainder relative to $\Sigma$ (lines 620, 649, 831).}
\emph{Before:} line~620 said that the uncertainty
$C_2\zeta^2\kappa_{\rm eff}^{-3}\le0.46\,\zeta^2\kappa_{\rm eff}^{-3}$ of the tail $\Sigma$ is
``$0.6\,\%$ of $\Sigma$ at $\kappa_{\rm eff}=29.6$''.  Lines~649 and~831 repeated this as a proved
relative error of at most $0.6\,\%$ at $\kappa_c\simeq30$.
\emph{Why this was wrong:} with $\Sigma\approx\zeta^2/(15\kappa_c^2)$ the ratio is
$C_2\,\zeta^2\kappa_c^{-3}\big/\bigl(\zeta^2/(15\kappa_c^2)\bigr)=15\,C_2/\kappa_c$.  With the
corrected $C_2\le0.47$ of C10 that is $0.238$ at $\kappa_c=29.6$, $0.235$ at $30$ and $0.219$ at
$32.2$; with the box coefficient $0.0641$ it is $0.47/(0.0641\cdot29.6)=0.248$.  With the old
$C_2\le0.46$ the figures were $0.233$ and $0.242$, so neither constant gives $0.6\,\%$.  The $C_1$ term adds less than $10^{-6}$ of $\Sigma$ there, even
at $\zeta=1/120$.  No computation in this document, in \texttt{numerics/pt\_checks.py},
\texttt{numerics/pt\_bars.py} or \texttt{rigor/pt\_tailconst.py}, in Note~5 or in
\texttt{certified\_numerics.tex} gives $0.6\,\%$ for this quantity.  The nearest figures describe
other quantities.  One is the measured deviation $\Sigma/(\zeta^2/(15\kappa_c^2))-1$ in the table of
\S\ref{sec:num}(4): $0.7$--$1.8\,\%$ at $\kappa_c=32.2$ and $2.2$--$6.3\,\%$ at $\kappa_c=18.4$, for
$\zeta=1/15,0.2,1$.  The other is the $C_2$ term as a fraction of $\Phi$, $0.2$--$0.3\,\%$ at
$\kappa_{\rm eff}=29.6$ for the tabulated $\zeta$.
\emph{Now:} the three lines give the bound, about $24\,\%$ ($25\,\%$ with the box coefficient
$0.0641$); lines 649 and 831 also quote the measured deviations.  The verdict after
Table~\ref{tab:bars} still says that Note~5's ``12--56\,\% off'' estimate was too pessimistic for the
occupation.  That now rests on the measured deviation, at most $6.3\,\%$ at $\kappa_c=18.4$ (the
cutoff of Note~5's estimate), and not on the bound, which is $15\cdot0.47/18.4=38\,\%$ there.  Source: found while making C1--C4; the check is the arithmetic above.

\paragraph{C6. A Note~5 table number (line 561).}
\emph{Before:} ``Note~5, Table~1 varies $(\Lambda,\kappa_{\max})$ over nine combinations at
$\zeta=1/15$''.  \emph{Now:} Table~2.  In the current Note~5 (\texttt{fidelity\_recovered\_quasifree.aux})
Table~1 is the large-$\zeta$ table \texttt{tab:largezeta}, and the nine-row convergence table
\texttt{tab:convergence} is Table~2.  This is the same renumbering as the Table~3$\to$4 of C3.

\paragraph{C7. The first-order remainder in overview item (B) (line 64).}
\emph{Before:} $|\tau_\pm-\nu_\pm-\zeta^2/(15\kappa_c^2)|\le C_2\zeta^2\kappa_c^{-3}
+C_1\zeta\kappa_ce^{-\kappa_c}$.  \emph{Now:} the last term is $C_1\zeta(1+\kappa_c)e^{-\kappa_c}$,
as Theorem~\ref{thm:tau} states and proves in \eqref{eq:taubound}; there $\tau_\pm-\nu_\pm=\Sigma$ by
\eqref{eq:phsym}, and $C_1=A_1/(4\pi^2)$ with $A_1$ defined through the factor
$(1+\kappa_c)^{-1}e^{\kappa_c}$.  With the same $C_1\le0.051$, the printed form was stronger than the
theorem on this term by the factor $(1+\kappa_c)/\kappa_c$, which is at most $1.125$ for
$\kappa_c\ge8$.  Knowledge-base card \texttt{tail-bound} states the $(1+\kappa_c)$ form.

\paragraph{C8. ``Certified'' error bars (lines 5, 75, 565, 589, 625, 631, 642, 826, 828).}
\emph{Before:} the title, the overview, \S\ref{sec:box}(b), the heading of \S\ref{sec:bars}, the
caption and header of Table~\ref{tab:bars}, the verdict after it and the summary called the bars of
this document ``certified''.  The phrases were ``certified error bars'', ``certified modulo the box
discretisation'', ``Certified bar'', ``certified interval'', ``The certified bar'', ``Certified
numerics'' and ``the certified upper bound''.
\emph{Why this was wrong:} every bar in this document is \eqref{eq:bar}.  Its upper end is
$\Phi^{\rm c}_{\rm sub}+0.173\,\zeta^2/\kappa_{\rm eff}+0.88\,\zeta^2/\kappa_{\rm eff}^2$, that is
\eqref{eq:final} at $\kappa_{\rm eff}$.  This is how \texttt{numerics/pt\_bars.py} computes it, and
the script reproduces the intervals of Table~\ref{tab:bars}, for example $[3.4866,\,6.5309]$ at
$\zeta=1/15$.  The constants of \eqref{eq:final} rest on the uncertified numerical inputs of C3 and
C4.  This comes on top of the box caveat of \S\ref{sec:box}(b) and the unaudited arithmetic margin
of \S\ref{sec:bars}.  None of the bars uses the bound $\Phi-\Phi_W\le\Phi_{22}+\Phi_{\rm off}$ of
\texttt{certified\_numerics.tex}, whose enclosures come from a different computation.  The lower
end, $\Phi^{\rm c}_{\rm sub}$, is a rigorous lower bound up to the arithmetic margin, as
\S\ref{sec:box}(a) says; that statement is unchanged.
\emph{Now:} these labels no longer say ``certified'', and the upper half is said to rest on the
numerical inputs of \eqref{eq:final}.  Three places are unchanged: ``not certified'' for the
extrapolated best value (line~629), and lines~608 and~644, which say what is still missing before
the digits can be called certified.  Source: knowledge-base card \texttt{tail-bound}, its Statement
(b) (explicit constants via numerical inputs) and its How to verify (the certified scripts
\texttt{numerics/certified/cert\_tail*.py} evaluate a different bound).

\paragraph{C9. A Note~5 corollary number (lines 73, 655, 712, 824).}
\emph{Before:} ``Note~5, Cor.~4.2'' in all four places.  \emph{Now:} Cor.~4.3.  All four cite the
corollary \texttt{cor:structure} of Note~5, ``Structure of the small-$\zeta$ expansion''.  Lines 73,
655 and 824 cite its $\log^{-2}(1/\zeta)$ rate, and line 712 its threshold $\kappa_*(\zeta)$.  The
same corollary is called Corollary~4.2 in \texttt{check\_note5\_theory.tex} (section
\texttt{sec:cor42}) and in \texttt{quadratic\_limit.tex} (line~26).  In the current Note~5 an
unlabelled remark (``Status of the proof'', after Proposition~4.1) takes the number 4.2, and
\texttt{fidelity\_recovered\_quasifree.aux} gives \texttt{cor:structure} the number 4.3.  This is the
same kind of renumbering as in C3 and C6.
The current version of that corollary also has a corrected statement (2026-09-09,
\texttt{rate\_of\_remainder.tex}): the expansion proceeds in integer powers, and ``the $\log^{-2}$
rate stated in an earlier version was wrong''.  Lines 73, 655 and 824 therefore now say that the
$\log^{-2}$ rate is the claim of an earlier version of that corollary.  Their comparison ``one
logarithm short of'' refers to that claim, not to the current statement.  Line 712 cites the
corollary for the threshold $\kappa_*(\zeta)$, which the current version still mentions, and is
changed only in the number.  Source: knowledge-base card \texttt{tail-bound} (``the $\zeta^2\log^{-2}$ remainder claimed
there was refuted'', card \texttt{rate-of-remainder}), and the check of the label above.

\paragraph{C10. The value of $A_3$, and hence of $C_2$ (lines 66, 328, 330, 332, 335, 445, 469, 470,
512, 517, 518, 522, 527, 531, 620, 649, 750, 751, 758, 764, 806, 812, 814, 831).}
\emph{Before:} \S\ref{sec:num}(3) read ``$\zeta^{-2}\|(g^{(2)})'''\|_{L^1}=0.0355,\,0.0326,\,0.0275,
\,0.0214$ at $\zeta=0.05,0.2,0.5,1$ (monotone decreasing in $\zeta$), so $A_3\le0.036$ and
$C_2=4\pi A_3\le0.46$''.  The document used $C_2\le0.46$, that is $4.6=10\,C_2$ in \eqref{eq:taunu},
throughout.
\emph{Why this was wrong:} $A_3$ is a supremum over $0<\zeta\le1$ (line~326).  If the values
decrease in $\zeta$, the supremum is the limit $\zeta\to0$, which was not sampled.  The functions of
\texttt{rigor/pt\_tailconst.py}, on its own grid ($t=0.05$ to $40$), give $0.03625$, $0.036406$ and
$0.036422$ at $\zeta=0.01$, $10^{-3}$, $10^{-4}$.  As $\zeta\to0$, $\zeta^{-2}g^{(2)}$ tends to the
closed form
\[
 g_0(t)=-\frac{1+\frac12\cosh t-\frac32\,t^{-1}\sinh t}{2\sinh^3(t/2)}=-\frac t{30}+O(t^3),
\]
because $\int_0^1\sinh^2\frac{ts}2\sinh^2\frac{t(1-s)}2\,\dd s=\frac14\bigl(1+\frac12\cosh t
-\frac32\,t^{-1}\sinh t\bigr)$.  Its $\|g_0'''\|_{L^1}$ is $0.03642$ on the script's grid.  On a
refined grid ($t=0.002$ to $80$, with $2\times10^4$ and $4\times10^4$ intervals, $h=10^{-4}$) it is
$0.03718$, and $0.03721$ when the part near $t=0$ is included.  At finite $\zeta$ the refined values
stay below this limit: $0.03698$, $0.03617$, $0.03331$ and $0.02175$ at $\zeta=0.01,0.05,0.2,1$.  So
$A_3=0.0372>0.036$ and $C_2=4\pi A_3=0.468>0.46$.  These are numerical values, not certified.
\emph{Now:} $A_3\le0.0374$ and $C_2\le0.47$, so $4.7=10\,C_2$ appears in \eqref{eq:taunu}, in the
bound on $\mathcal T$ of Theorem~\ref{thm:main} and in $s(\kappa_c)$.  The numbers that depend on
$C_2$ were re-derived.  $s(10)=0.667+0.47+0.494+0.077=1.707\le1.75$ (C2), so $\mathcal
T\le1.75\,\zeta^2/\kappa_c^2$, $\mathcal T\le0.0175$ and $Y\le0.018$ are unchanged.  The corrected
chain against \eqref{eq:final} has supremum $0.98405$ at the corner ($0.98467$ with $f=0.008516$), so
\eqref{eq:final} still holds given the inputs (C3).  The $C_2$ term is about $24\,\%$ of $\Sigma$
at $\kappa_{\rm eff}\simeq30$, or $25\,\%$ with the box coefficient (C5).  The factor to spare at
line~764 becomes $6$--$31$.  Scripts: \texttt{numerics/kb-checks-2026-09-24/a3\_refined.py}, which
needs mpmath, and \texttt{refb10\_tail\_sup\_c2.py}, a copy of \texttt{refb10\_tail\_sup.py} with
$4.7$ and a scan of the whole grid.  Each has its \texttt{.out}, and the folder has a
\texttt{README.md}.  Source: referee report REF-DOC-TAIL (2026-09-24), checked independently by the
fix author with \texttt{a3\_refined.py}.  Knowledge-base card \texttt{tail-constants-numerical} still
records $A_3\le0.036$.
\end{document}
